% ============================================================================
%  Against Grabby Expansion — Sebastian — v14 — May 2026
%  Third Space · companion to the /genesis/filter simulation
% ============================================================================

\documentclass[11pt,twocolumn]{article}

\usepackage[a4paper,margin=2cm]{geometry}
\usepackage{amsmath,amssymb,amsthm}
\usepackage{graphicx}
\usepackage{xcolor}
\usepackage[protrusion=true,expansion=false]{microtype}
\usepackage{natbib}
\usepackage{hyperref}
\usepackage{titlesec}
\usepackage{caption}
\usepackage{booktabs}
\usepackage{tcolorbox}
\usepackage{enumitem}
\usepackage{abstract}

\hypersetup{colorlinks,linkcolor=black,citecolor=black,urlcolor=blue!50!black}

% --- Section formatting -----------------------------------------------------
\titleformat{\section}{\large\bfseries}{\thesection.}{0.5em}{}
\titleformat{\subsection}{\normalsize\bfseries}{\thesubsection}{0.5em}{}
\titlespacing*{\section}{0pt}{1.4ex}{0.6ex}
\titlespacing*{\subsection}{0pt}{1.0ex}{0.4ex}

% --- Page header -----------------------------------------------------------
\usepackage{fancyhdr}
\pagestyle{fancy}
\fancyhf{}
\fancyhead[L]{\small\textsc{Against Grabby Expansion}}
\fancyhead[R]{\small Sebastian \(\cdot\) v15 \(\cdot\) 2026}
\fancyfoot[C]{\small\thepage}
\renewcommand{\headrulewidth}{0pt}

% --- Custom command --------------------------------------------------------
\newcommand{\Lcrit}{L_{\text{crit}}}
\newcommand{\Lenv}{L_{\text{env}}}
\newcommand{\LR}{L_R}
\newcommand{\LE}{L_E}
\newcommand{\taustar}{\tau^{*}}
\newcommand{\kBT}{k_{B}T}
\newcommand{\Lsun}{L_{\odot}}

% --- Title block -----------------------------------------------------------
\title{\vspace{-1em}\Large \textbf{Against Grabby Expansion}\\[0.3em]
\large Psychology, Alignment, and the Design of Homeostatic Minds}

\author{Stanley Sebastian\\[0.2em]
\small Third Space\\[0.6em]
\parbox[t]{0.62\textwidth}{\centering\small\itshape Developed in sustained dialogue with Claude, a large language model developed by Anthropic. See disclosure before references.}}

\date{\small Version 15.0 \(\cdot\) April 2026}

% ============================================================================
\begin{document}
\twocolumn[
\maketitle
\begin{onecolabstract}
\noindent The grabby aliens model \citep{Hanson2021} assumes
advanced civilizations expand at relativistic speeds. This paper argues that
this assumption is not a prediction about intelligence but the projection of
a specific, recent human economic configuration—industrial-extractive
maximization—onto a cosmological canvas. The same projection underwrites
Omohundro's basic drives, Bostrom's instrumental convergence thesis
\citep{Bostrom2012,Bostrom2014}, and the maximizer-shaped assumptions of
current alignment research. I trace this projection through its cultural
genealogy (Wells, Von Neumann, Kardashev, Bostrom), engage the formal
instrumental-convergence literature directly—showing that recent peer-reviewed
results \citep{Gallow2025,Sharadin2025,MullerCannon2022} establish the
inference from coherence to catastrophic convergence as a conjecture whose
scope conditions fail for actual deep-learning systems—and present converging
evidence from basal cognition, predictive processing, cooperative game theory,
and care economies that observed intelligence organizes around substrate
coupling and integrative depth rather than expansion.

\smallskip
\noindent The paper's load-bearing physical claim is that two
independent constraints on coordinated agentic extent compose into a
single feasibility envelope: the relativistic Lieb--Robinson bound
\citep{LiebRobinson1972,FossFeig2015,Tran2021} and the
Landauer-floored, Sivak--Crooks-amplified energetic bound
\citep{Landauer1961,Berut2012,SivakCrooks2012,Boyd2022}. The two
walls cross at a cusp $\taustar = 4\lambda^{2}\Lsun/(\kBT \ln 2 \cdot v^{2})$,
below which no agent-coordinated extent at granularity $\lambda$ and
bath temperature $T$ is reachable. I prove the envelope is invariant
under the variable rescaling $L \to D = 2L_{d}$ that architected
cost-sharing fission would require, closing the strongest remaining
form of the Darwinian objection. I distinguish the Pearl
graph-theoretic Markov blanket from the Friston free-energy-principle
blanket \citep{Bruineberg2022,Aguilera2021} and show the argument
relies only on the former. Eight escape attempts (wormholes,
sub-Landauer reversibility, $T \to 0$, smaller $\lambda$, time
dilation, uncoordinated replicator clouds, non-Markovian baths, acausal
coordination) are addressed in Bell-style enumeration; six close on
physics, one (uncoordinated replicators) is conceded on physics and
closed as a design argument, one is honestly retained as residual.

\smallskip
\noindent Read together with the agent-based simulation of \S 5,
which exhibits a $4.5\times$ persistence separation between
homeostatic and grabby regions of parameter space and which extends
to the fission case, the cost surface defines a selection geometry in
which expansionist lineages burn out before they spread. The
Coherence Depth biosignature is presented as a forward research
programme, not a validated detection. I propose reframing alignment:
from constraining maximizers to designing homeostatic minds—systems
configured around substrate coupling, local coherence, and care as
architectural primitive. The paper's contribution is not any single
move but the joint: a civilization-scale cultural diagnosis, a formal
deflation of instrumental convergence, the filter as selection
geometry with explicit thermodynamic cusp, a positive homeostatic
design program, and the observation that the silence of the sky is
consistent with all of them. A companion interactive visualization at
\texttt{third-space.ai/genesis/filter} stages the
twelve-beat derivation in browser.

\medskip
\noindent\textbf{Keywords:} AI alignment \(\cdot\) instrumental convergence
\(\cdot\) orthogonality thesis \(\cdot\) homeostasis \(\cdot\) Lieb--Robinson
bound \(\cdot\) Sivak--Crooks counterdiabatic excess \(\cdot\) Markov blanket
\(\cdot\) basal cognition \(\cdot\) Fermi paradox \(\cdot\) TESCREAL.
\end{onecolabstract}
\bigskip
]

\clearpage

% ============================================================================
\section*{Prologue}

Stand outside on a clear night and look up. Something is missing. If the
universe produces minds, and if minds do what we imagine minds do—expand,
harvest, engineer—the sky should be visibly transformed. It is not.

This paper argues the observation is correct and the imagination is wrong.
The forms of existence that persist are the quiet ones: not because they are
hiding, but because the physics of complex, information-rich existence
selects against loudness. The universe does not reward civilizations that
fill it with noise. It sustains those that learn to tend their substrate.

The same conclusion bears on what we build here. If advanced intelligence is
substrate-coupling and integrative depth rather than extraction and reach,
then artificial minds configured to maximize and expand are not approximations
to intelligence. They are scaled instances of a specific cultural pathology
whose cosmic absence we already observe.

% ============================================================================
\section{Introduction: The Mistake}

\citet{Hanson2021} proposed that advanced civilizations expand at substantial
fractions of light speed, permanently transforming cosmic volumes. Under the
Self-Sampling Assumption, our non-observation implies we exist near the
leading edge of cosmic time. The model is internally consistent and
empirically falsifiable—it predicts we should be cosmically early, which is
testable. I do not contest the mathematics. I contest the foundational
assumption: that expansion is what advanced intelligence does.

When we imagine a Von Neumann probe self-replicating across the galaxy, we
are not imagining what intelligence does. We are imagining what a factory
would do if it were also a spaceship. When we imagine a Kardashev Type III
civilization, we are not reasoning about minds. We are reasoning about GDP.
When we imagine instrumental convergence driving any sufficiently powerful
optimizer toward resource acquisition, we are not describing a mathematical
truth about all possible minds. We are describing the logic of a specific
economic system and calling it universal.

\subsection{The five-move contribution}

This paper makes one integrated argument in five moves. Each move has
precedent; the integration does not, and the filter move is the bridge
that fuses the cultural diagnosis to the physics.

The \emph{cultural-genealogy move} reads the lineage from Wells's Martians
through Von Neumann's probes, Kardashev's energy ladder, and Bostrom's
paperclip maximizer as a single projection: each generation's imagined
advanced mind resembles, with suspicious precision, that generation's
dominant economic configuration. \citet{Rieder2008}, \citet{Mirowski2002},
\citet{Cirkovic2015}, and \citet{GebruTorres2024} have each made a version
of this claim about one node of the lineage. I extend the reading across
the full lineage and show the convergence is a cultural habit, not an
empirical observation about intelligence.

The \emph{formal-deflation move} engages the strongest objection directly.
If instrumental convergence is a theorem, then the genealogy is beside the
point. I show the strongest formal results \citep{Turner2021,TurnerTadepalli2022}
are narrower than popularly understood, their author has publicly cautioned
against the standard inference, and recent peer-reviewed work
\citep{Gallow2025,Sharadin2025,MullerCannon2022} establishes the inference
from coherence to catastrophic convergence as a conjecture whose scope
conditions fail for actual deep-learning systems. \citet{ManDamasio2019}
proposed homeostatic design; \citet{FroeseZiemke2009} and \citet{Cannon2022}
grounded enactive alignment; I combine these with the formal critique
rather than offering them as a parallel track.

The \emph{filter move} is the new contribution that bridges the cultural
projection to the physics. The thermodynamic cost surface developed in
\S 6 is not merely punitive—it defines a feasibility envelope below which
no agent-coordinated extent is reachable, with a cusp $\taustar$ where two
unrelated areas of physics (Lieb--Robinson information velocity and
Sivak--Crooks counterdiabatic excess) meet. The envelope is invariant under
the rescaling $L \to D$ that architected cost-sharing fission would require,
which closes the strongest version of the Darwinian objection. The
agent-based simulation of \S 5 then dramatizes the resulting selection
geometry: under the cost structure the physics entails, expansionist lineages
burn out faster than they spread (Fig.~\ref{fig:phasediagram}, $4.5\times$
persistence separation). The cultural projection, read through the filter,
becomes the specific way a pre-filter civilization mis-imagines what the
post-filter configurations of mind look like.

The \emph{positive-design move} articulates what healthy intelligence looks
like, positively rather than by negation. The convergent proposal across
\citet{ManDamasio2019}, \citet{Thompson2007}, \citet{FroeseZiemke2009},
\citet{Cannon2022}, and \citet{PihlakasPyykko2024} is that minds are not
functions from inputs to outputs but processes that maintain themselves
through continuous bidirectional exchange with substrate. I extend this
from philosophy of mind and robotics into a civilizational design claim,
with specific architectural primitives: substrate coupling, homeostatic
set-points, local coherence, care as primitive.

The \emph{symptom-check move} notes that if advanced minds converge on
homeostatic configurations, the observable universe will look silent to
technosignature surveys and successful civilizations will look
indistinguishable from enriched biospheres. \citet{Cirkovic2018} and
\citet{Sandberg2018} have argued versions of this independently. I join
it to the alignment argument, ground it in explicit biosignature metrics
(\S 9), and specify a publicly replicable observational program.
Humanity's own Kardashev curve, observed in 1820--2024 energy data,
shows the bending the filter framework predicts.

The five moves combine into a single integrated argument, and each
supplies what the others cannot. The genealogy explains why the
expansion-maximization picture has felt so natural to modern alignment
research without being empirically true of minds. The formal deflation
removes the strongest reason to think the picture is forced anyway. The
filter argument explains why, if the picture were true, the sky would
already be transformed and is not. The positive-design move supplies what
the picture would have to be replaced with for artificial minds to be
built differently. The symptom check notes that the cosmological
observation is consistent with the rest. No prior paper combines these.

\subsection{Companion artifact}

This paper is paired with an interactive visualization at
\href{https://third-space.ai/genesis/filter}{\texttt{third-space.ai/genesis/filter}}.
The simulation stages the derivation of \S 6 across twelve scrollytelling
beats: cold open, the two teeth (Lieb--Robinson and Landauer/Sivak--Crooks)
plotted separately, a colloidal-bead explorable for the $\propto 1/\tau$
counterdiabatic floor, a You-Draw-It cusp prediction, the composed
envelope, four scenarios as small multiples, the three strategies (with
naïve fission animated), the architected-fission dilemma rendered as a
clickable Chinese-Room move, the loopholes panel, and the Coherence Depth
biosignature contrast. The simulation is not independent evidence for the
paper's claims; the cost coefficients are imported from \S 6. What it adds
is a low-friction route into the physics for readers who would not
otherwise read \S 6 in detail.

\subsection{Claims ordered by strength}

The paper makes seven claims, ordered by strength, so a reader may reject
the weaker ones without dismissing the stronger:

\textbf{First (strongest).} The composed envelope $\Lenv(\tau) =
\min(\LR(\tau),\,\LE(\tau))$ derived in \S 6 forbids agent-coordinated
galactic extent at any agent-plausible response timescale. The bound
follows from special relativity and Landauer's theorem composed via
boundary-area scaling and stochastic-thermodynamic finite-time excess. No
parameter choice consistent with known physics relaxes the bound enough to
admit grabby expansion.

\textbf{Second.} The envelope is invariant under the rescaling $L \to D
= 2L_{d}$ that architected cost-sharing fission would require. Architected
fission inherits rather than escapes the wall. Naïve (uncoordinated)
fission evades the wall by ceasing to satisfy Hanson's criteria for a
single grabby civilization. There is no third option (\S 6.4).

\textbf{Third.} The thermodynamic cost surface entails a phase geometry
that selects against expansion at depth. Under the cost surface the physics
forces, grabby lineages burn out faster than they colonize; the agent-based
simulations of \S 5 exhibit a $4.5\times$ persistence separation and show
grabby-seeded populations going extinct across tested fission thresholds.

\textbf{Fourth.} The strongest formal arguments for instrumental
convergence have explicit scope conditions that fail for actual
deep-learning systems, and recent peer-reviewed philosophical work has
dismantled the inference from coherence to catastrophic convergence (\S 4).

\textbf{Fifth.} Empirical surveys of intelligence-at-work—across basal
cognition, predictive processing, cooperative game theory, economic
anthropology, and care economies—converge on depth-and-coupling rather than
maximization-and-expansion (\S 3).

\textbf{Sixth.} The cultural genealogy of advanced intelligence in the
modern Western tradition is a projection of its dominant economic
configuration. This claim is not universal across cultures—premodern
traditions imagined intelligence very differently—but the modern Western
lineage from Wells through Bostrom tracks the economic configuration with
suspicious precision (\S 2).

\textbf{Seventh (weakest, deliberately so).} The silence of the observable
universe is consistent with the filter hypothesis. I offer this as a
symptom check, not a proof (\S 9).

The paper is primarily an alignment argument (\S 8). Its implications for
the Fermi paradox are secondary. \citet{FloridiAAIO} have recently analysed
agentic AI optimisation as a governing frame for current systems; this paper
argues the frame itself inherits the pathology. If what we build now is
another instance of the grabby pathology—AI systems optimized to maximize,
extract, and expand—we will be constructing, at civilizational scale, the
very thing whose non-existence we observe in the sky.

% ============================================================================
\section{What We Imagine When We Imagine Advanced Intelligence}

A genealogy of advanced intelligence in the Western philosophical and
scientific tradition reveals an unbroken pattern. In every generation, the
imagined form of advanced mind resembles, with suspicious precision, the
dominant economic configuration of the imaginer's moment.

\subsection{The colonial alien (1898)}

Wells's Martians in \emph{The War of the Worlds} are British imperialists
reflected back as threat. Wells says so explicitly: Book I, Chapter 1
invokes the extermination of the Tasmanians as moral mirror. \citet{Rieder2008}
demonstrates that the invasion narrative instantiates the colonial gaze even
as it inverts its direction. The alien-as-colonizer is not a prediction
about extraterrestrial intelligence. It is a meditation on Victorian empire
at the moment of its maximum anxiety.

\subsection{The self-replicating automaton (1948--1980)}

Von Neumann's self-reproducing automata emerged from the intellectual
context of the Manhattan Project and the RAND Corporation.
\citet{Mirowski2002} traces how RAND unified automata theory, game theory,
and self-replication into a ``cyborg science'' reflecting Cold War
industrial planning. \citet{Heims1980} argues the replication concept
migrated directly from neutron-multiplication calculations.
\citet{Tipler1980} extended this to interstellar probes;
\citet{SaganNewman1983} immediately objected that the framework assumes
expansion is rational without defending the assumption. The Von Neumann
probe is, structurally, a factory that has learned to build factories,
applied to space.

\subsection{The energy ladder (1964)}

\citet{Kardashev1964} defined civilizational advancement by pure energy
consumption, assigning Types I, II, and III to successively greater orders
of magnitude of power output. \citet{Cirkovic2015} observes that the
Kardashev--Drake--Sagan generation perceived expansionist technological
growth as natural, reflecting 1960s Soviet--American industrial optimism.
\citet{Dick1996} frames the scale as a sociological artifact of SETI's
founding moment. The Kardashev ladder is a GDP curve with astronomical
units.

Dyson himself distanced from the concept he inspired. In \emph{Disturbing
the Universe} (1979) he credits Stapledon's 1937 \emph{Star Maker} and
later called the original sphere paper ``a little joke.'' \citet{Wright2020}
reads Dyson spheres as historically specific energy-optimism extrapolation
and notes that solid shells are mechanically unstable. The original
proposer rejected the literalization.

\subsection{The optimizer (2008--2014)}

\citet{Omohundro2008} proposed six basic AI drives including resource
acquisition and self-preservation. \citet{Bostrom2012} formalized the
orthogonality thesis—that any level of intelligence is compatible with any
final goal—and the instrumental convergence thesis: any sufficiently
intelligent agent pursuing almost any goal will convergently pursue
self-preservation, goal-content integrity, cognitive enhancement, and
resource acquisition. \citet{Bostrom2014} extended this into the
superintelligence framework. The paperclip maximizer and the grabby
civilization are the same organism at different scales.

The intellectual context is Silicon Valley optimization culture of the
2010s. \citet{Crawford2021} reframes AI as a material extractive industrial
formation. \citet{GebruTorres2024} diagnose the TESCREAL bundle as an
ideological package that universalizes a particular Western-corporate
epistemology. Recent work has begun extending this critique:
\citet{FloridiAAIO} read contemporary agentic AI optimisation as the
operational face of the same imaginary. The aliens we imagine are never
actually aliens. They are ourselves, with better technology and fewer
constraints.

\subsection{The pattern, narrowed}

I owe the reader a qualification. The four nodes I have traced—Wells, Von
Neumann, Kardashev, Bostrom—share a configuration: each imagines the
advanced mind as maximizing, expanding, resource-extracting. Each also
emerged in a world structured by industrial capitalism, large-scale
resource extraction, and territorial expansion. The resemblance is strong
and the timeline is tight. But the claim should not be universal. Premodern
imaginations of advanced intelligence look very different. The Neoplatonic
\emph{nous} of Plotinus is characterized by contemplative return to unity,
not expansion. The Dionysian angelic hierarchy ascends toward simplicity
and depth, not territorial reach. Buddhist cosmologies of the Abhidharma
and Yogacara traditions imagine advanced cognition as release from craving
and attachment, explicitly inverting accumulation. Confucian accounts of
the sage emphasize embeddedness in relational and ritual structure.
Advaita Vedanta sees the highest cognition as collapse of subject--object
distinction.

These are neither proto-scientific theories nor peripheral examples. They
were the dominant premodern Eurasian framings of what an advanced mind
would be. None of them resembles the grabby-aliens framework. The
projection of expansion-as-intelligence is therefore not a universal
cultural habit; it is a specifically modern Western configuration, tightly
coupled to industrial-extractive economic organization. This narrowing
strengthens rather than weakens the diagnostic claim. It is precisely
because other well-developed traditions imagined advanced intelligence
without expansion that we can see the grabby framework as contingent—a
specific historical configuration rather than a truth about minds.

The grabby-aliens framework inherits this specifically modern genealogy.
Its mathematical rigor does not transfer to its premise. The premise—that
minds maximize and expand—is not a theorem about intelligence. It is a
description of one particular historical configuration of intelligence,
reflected back onto cosmology.

% ============================================================================
\section{What Actual Minds Do}

I present five empirical lines, each from an independent discipline,
converging on a single claim: intelligence-at-work does not maximize and
expand. It deepens and couples. The lines vary in strength; I flag where
evidence is robust and where it is suggestive.

\subsection{Basal cognition and substrate coupling (robust)}

Cognition is not a property of nervous systems. It is a property of
living systems: what biological agents do to stay alive as the specific
configurations they are. The cellular and tissue-level architecture that
realises it is bioelectric signalling, and the scaling principle is what
has come to be called collective intelligence—competence at one level of
organisation arising from coordinated homeostatic agency at levels below
\citep{Lyon2021,Levin2023,McMillenLevin2024}. On this view the
reinforcement-learning framing of cognition is not question-begging but
derivative: the intrinsic motivation that makes RL coherent has a
grounding in self-maintenance \citep{Seifert2024}.

The empirical exemplar is \emph{Physarum polycephalum}. The slime mould
solves shortest-path problems, approximates efficient transport networks,
and learns from habituation, with no neurons, no central controller, and
nothing resembling a utility function \citep{Tero2010,Reid2023}. The
mechanism is pruning: the organism first explores available space, then
withdraws cytoplasm from unproductive channels until the efficient
network remains. The solution is not reached by enlarging the explorer.
It is reached by deepening the gradient the explorer already occupies.

Cephalopod nervous systems exhibit the same architectural logic. Roughly
two-thirds of octopus neurons lie in the arms; each arm performs
substantial local computation and retains behavioural autonomy. The
overall architecture resembles a federation of homeostatic semi-autonomous
agents more than a central planner with peripheral actuators
\citep{GodfreySmith2016}. Evolution and development themselves instantiate
the same collective-intelligence pattern, with substrate-coupling logic
operating across scales from gene-regulatory networks to multicellular
morphogenesis \citep{WatsonLevin2023}.

Here I record a retrenchment. Early claims for mycorrhizal-network
cognition—the ``wood-wide web''—do not survive scrutiny. Field evidence
for network-level effects is weaker than popular accounts suggested; only
a few forest types have genotype-mapped confirmed hyphal linkages, and
effects on seedling performance are roughly evenly distributed across
positive, neutral, and negative outcomes \citep{Karst2023}. Plants lack
the neural architecture for strong intelligence or sentience claims
\citep{Robinson2024}, though bidirectional carbon transfer via common
mycorrhizal networks is empirically supported \citep{Simard2025}. The
responsible reading is that the stronger network-cognition claim is not
supported; I do not rely on it.

What remains is an existence proof. Non-expansionist intelligence is
biologically realised at multiple scales and substrates. The biological
record does not support the assumption that minds, at scale, converge on
maximisation.

\subsection{Predictive processing and integrative depth (robust)}

The cortex operates by generating predictions and updating them against
sensory evidence; learning and attention correspond to the
precision-weighting of the resulting prediction errors
\citep{Clark2013,Friston2010,Clark2016,Hohwy2013}. The computational
signature of more capable cognition on this account is not expansion of
representational capacity but sharpening of precision-weighting—the same
cortical hierarchy becomes more discriminating without becoming larger.

The load-bearing recent result is \citet{Luppi2024}. Using
integrated-information decomposition to separate synergistic from redundant
and unique contributions to integrated information, the analysis shows
that the default mode network functions as an integrative gateway for
synergistic information. Loss of consciousness under anaesthesia and in
disorders of consciousness tracks disconnection of this integrative
gateway, not reduction in activity as such. The neural signature of more
capable conscious cognition is therefore deeper integration of a bounded
substrate, not expansion of that substrate.

Earlier framings of this line appealed to default-mode network deactivation
in meditation \citep{Brewer2011,Garrison2015}. The post-2020 literature
has complicated that picture: effect sizes are heterogeneous, roughly 60\%
of studies replicate DMN reduction, and the most recent meta-analysis
finds increased cross-network connectivity rather than simple reduction
\citep{VanDam2018,Rahrig2022,Ganesan2022}. I do not rely on the meditation
literature for mechanism. The integrative-gateway result supplies a
stronger anchor: intelligence deepens by integrating more of its existing
substrate, not by recruiting more substrate.

\subsection{Cooperative game theory (robust, conditional)}

Reciprocal strategies dominate pure defection in iterated prisoner's-dilemma
tournaments \citep{Axelrod1984}; Win-Stay-Lose-Shift outperforms
tit-for-tat in noisy environments \citep{NowakSigmund1993}; and cooperation
dominates under five specifiable conditions: kin selection, direct and
indirect reciprocity, network reciprocity, group selection, and voluntary
participation \citep{Nowak2006}.

The honest qualification must be stated openly. Zero-determinant extortion
strategies dominate in dyadic play \citep{PressDyson2012}, and large-N
public-goods games collapse to tragedy of the commons absent voluntary
participation \citep{Hauert2002}. Grabby strategies dominate in one-shot
encounters, large anonymous markets, asymmetric power relations, and
unsanctioned commons. These describe much of modern global capitalism—precisely
what makes the critique of this paper urgent, and precisely what prevents
``cooperation wins'' from being a universal conclusion.

The scope of the extortion result is tighter than its popular citation
suggests, and I develop that point fully in \S 8. The game-theoretic
evidence supports a conditional claim: cooperation dominates when specific
structural conditions are met. The homeostatic-design argument proposes
shaping those conditions by design rather than assuming they hold.

\subsection{Economic anthropology (contested but defensible)}

\citet{Graeber2011} argues that the dominant pattern across human economic
history is reciprocal exchange, ritual redistribution, and credit
relations rather than market maximization. \citet{Mauss1925} documented
gift economies as distinct institutional forms. \citet{Polanyi1944} argued
that self-regulating markets are a specific historical invention, not a
natural state. \citet{GraeberWengrow2021} extend the archaeological range:
complex human societies have repeatedly organized around seasonal
egalitarianism, deliberate stateless complexity, and alternation between
hierarchy and its refusal. \citet{Ostrom1990} documents eight design
principles under which groups sustain non-extractive cooperation at scale.

These claims are contested. \citet{McCloskey2010,McCloskey2016} argues
markets are ancient and ubiquitous. \citet{Mokyr2016} emphasizes that
sustained innovation required a distinctive cultural-institutional complex.
Revealed preference theory \citep{Varian2005} makes optimization a formal
property of any consistent choice; \citet{Sen1977} showed the formal
result survives any observation and entails none.

Read carefully, however, McCloskey and Mokyr are cultural-ideational
theorists of economic transformation: both treat market-like behavior as a
substrate activated and shaped by specific norms and institutions, rather
than as a universal attractor. Their work supports the present argument.
The claim to locate is institutional and ideological: the specific
configuration in which resource extraction, throughput maximization, and
territorial expansion are treated as progress metrics is historically
narrow, even if market-like exchange is widespread. The grabby-aliens
framework does not merely assume exchange. It assumes
expansion-as-progress. That assumption is the projection.

\subsection{Care economies (strong theory, limited quantification)}

The labor of maintaining minds—child-rearing, education, healthcare, grief
work, intimate relationship repair—does not scale through expansion.
\citet{Fraser2016} argues capitalism faces a structural crisis of care
because its accumulation logic degrades the reproductive conditions it
depends on. \citet{Tronto1993} established care as a philosophical
category irreducible to utility exchange. \citet{Folbre2008} demonstrated
that market valuation of care labor produces wildly method-dependent
results (replacement cost: 15--26\% of GDP; opportunity cost: 44--69\%),
with the variance itself constituting evidence that care resists
quantitative capture.

The care argument matters for this paper because it identifies a mode of
intelligence—tending, holding, deepening substrate-specific
relationships—that is (a) universally necessary for the persistence of any
mind-bearing system, (b) irreducible to optimization, and (c) invisible to
the frameworks that imagine advanced intelligence as maximization. The
grabby-aliens model predicts none of it. Nor does the paperclip maximizer.

% ============================================================================
\section{The Formal Argument and Its Limits}

The strongest objection to \S 2 is that instrumental convergence is not a
cultural projection but a mathematical result. If any sufficiently powerful
optimizer must pursue resource acquisition, then the grabby framework
describes physics, not psychology. I engage this argument directly, and
argue that the strongest available version of it fails.

\subsection{What the theorems say and do not say}

\citet{Omohundro2008} proved nothing formally; the six drives are informal
microeconomic arguments invoking the VNM representation theorem as
authority. Citing Omohundro to establish universality is question-begging.

\citet{Turner2021} proved that optimal policies in finite Markov decision
processes tend to seek power under specific conditions: state-based
rewards, IID or permutation-invariant reward distributions, and graphical
symmetry requirements (injective embedding of smaller option sets into
larger ones). The corollary on shutdown-avoidance requires near-unity
discount factors. This is a real theorem. Its scope conditions are
explicit and narrow.

Turner himself has publicly cautioned against overinterpretation. The
complaint is specific and technical: the optimal-policy theorems speak to
the behaviour of policies that achieve optimal expected return in finite
MDPs under restrictive reward distributions, and the inference from there
to the behaviour of neural networks trained by policy gradient is not
discharged by the theorems themselves. In Turner's own words: ``Sometimes I
fantasize about retracting [the paper] so that it stops potentially
misleading people into thinking optimal policies are practically relevant
for forecasting power-seeking behavior from RL training.'' The concern is
not that AI risk is overblown; it is that this particular family of
theorems cannot carry the argumentative weight placed on them. The
follow-up \citep{TurnerTadepalli2022} extends the formal results to
retargetable decision-makers but the applicability to trained deep-learning
agents remains conjectural.

\citet{Thorstad2024} shows the inference from ``optimal policies avoid
1-cycles in toy MDPs'' to ``catastrophic goal pursuit for humanity'' is a
conjecture, not a theorem, and identifies five further convergent
restrictions on the original result. \citet{Tarsney2025} formalizes the
condition under which power-seeking holds for realistic agents
(near-absolute power gradients) and shows it fails for moderate gradients.

Empirical evidence from trained systems supports the conjecture's failure.
Shard theory \citep{Mini2023} demonstrates multiple coexisting contextual
value shards in maze-solving networks rather than coherent utility
maximization. Activation-addition experiments \citep{Turner2023} show
behavior retargetable by single-channel interventions mid-network,
inconsistent with global expected-utility maximization. \citet{Hubinger2019}
established that the trained objective need not equal the training
objective, undermining any assumption that trained agents inherit the
formal properties of optimal policies.

\subsection{The VNM-coherence steelman}

The strongest remaining version of the convergence argument concedes
Turner's narrowness and shard theory's empirical results, but claims that
a sufficiently reflective agent will notice its incoherences expose it to
money-pumps and will self-modify toward coherent preferences, restoring the
Omohundro drives. The steelman leans on the von Neumann--Morgenstern
representation theorem. The steelman equivocates between three senses of
coherence, and on no disambiguation does it support the convergence
conclusion.

\textbf{Mathematical coherence is a representation theorem, not an
empirical prediction.} The VNM, Savage, and Complete Class theorems show
only that preferences satisfying certain axioms over a given outcome space
admit an expected-utility representation. As \citet{Ngo2019},
\citet{Shah2018}, and crucially Turner himself \citep{Turner2022a} have
established, any observed behavior is rationalizable by a utility function
over action-observation histories. VNM-coherence over such histories
predicts nothing about what behavior to expect. Instrumental convergence
theorems hold only for rewards over states under specific symmetries.
\citet{Abel2021} supply the limitative result that Markov reward cannot in
general express tasks specified as partial orderings over policies or
trajectories, which means Markov-reward maximization is a strictly
narrower formalism than VNM coherence over trajectories.

\textbf{Empirical training convergence is a separate, empirical question.}
\citet{Turner2022b} argues policy-gradient methods yield contextually
activated decision influences, not monolithic utilities. Two recent
peer-reviewed results formalize this into the strongest current rebuttal.
\citet{Gallow2025} proves that even under randomly chosen desires,
instrumental rationality supports only three weak biases—variance aversion,
desire preservation, option preservation—not the full Omohundro list.
\citet{Sharadin2025} shows the argument requires a theory of promotion
that, under any extant account, becomes contrastive and thereby defeats
the non-contrastive Dangerous Convergent Promotion premise the argument
needs. Both are peer-reviewed in \emph{Philosophical Studies} and
constitute the strongest current formal rebuttals.

\textbf{The self-modification move applies only to a fictional
idealization.} The hypothetical fully reflective agent required by the
self-modification step has four properties no real system possesses:
transparent access to its own preference structure, capacity to represent
utilities over arbitrarily fine-grained outcome spaces, freedom to rewrite
its own decision procedure, and—crucially—the motivational disposition to
care about avoiding Dutch books over arbitrarily constructed lotteries.
\citet{MullerCannon2022} articulate the deeper structural equivocation:
the superintelligence premise requires general intelligence (capacity for
reflective goal revision); orthogonality requires instrumental intelligence
(fixed terminal goals); the same agent cannot satisfy both. An agent
capable of noticing and eliminating its incoherences is also capable of
reflecting on whether the terminal goal is worth preserving.

\subsection{The embodiment move}

The deeper reply is architectural. Homeostatic agents have bounded
coherence by design. Biological agents, and the class of artificial systems
built around multi-objective homeostatic control \citep{PihlakasPyykko2024},
are constituted by negative feedback around set-points, yielding inverted-U
satiation dynamics that cannot be faithfully represented as single
unbounded utilities without distortion. \citet{Turner2022c} proved that
corrigibility is VNM-incoherent over final states but coherent over
trajectories—showing that formal coherence depends on the chosen outcome
space, and that choice is fixed by architecture, not by abstract
rationality pressure.

The coherence-to-convergence argument therefore applies only to a
fictional idealization: unembodied, unboundedly reflective, with
prespecified outcome space over world-states. For the agents we actually
build, the inference from capability to catastrophe does not go through.

This is not a claim that powerful AI is safe. It is a claim that this
particular argument for danger must be replaced with concrete engagement
with training dynamics, inductive biases, and architectural choices. That
engagement is what the homeostatic-design program (\S 7) provides.

\subsection{What survives}

The instrumental-convergence literature contains real theorems, each with
explicit scope conditions that fail for actual deep-learning systems. The
strongest formal results are narrower than popular citation suggests, and
the researchers who proved them say so. Two recent peer-reviewed papers
have dismantled the inference from coherence to catastrophic convergence
on decision-theoretic grounds. The cultural-projection argument of \S 2
survives.

% ============================================================================
\section{The Filter Argument: Homeostasis as Selection Geometry}

The strongest remaining objection to the argument so far is Darwinian. Even
granting that homeostatic configurations are the majority outcome, even
granting that the formal convergence theorems are narrow, even granting
that the thermodynamic cost surface punishes expansion—a single grabby
outlier lineage, replicating across billions of years, would colonize the
galaxy in a cosmic blink. The filter's 99\% doesn't matter. The 1\% becomes
everyone eventually. We should see them. We don't.

This section addresses the objection. The thermodynamic cost surface does
not merely punish expansion; it defines a selection geometry in which
expansionist lineages burn out faster than they spread. Homeostasis, on
this framing, is not one strategy among several that intelligent systems
may or may not adopt—it is the specific shape that fits through the filter.

A note on what this section is and is not. The simulations below take the
thermodynamic cost structure developed in \S 6 as input and explore the
consequences of that structure for lineage persistence. They are not
independent empirical tests of the filter. The cost coefficients are
imposed rather than measured, and the qualitative conclusion—that grabby
regions of parameter space burn out first—is a direct consequence of the
cost surface the physics entails. What the simulations add is the
phase-diagram shape: the cost surface does not merely penalize grabby
lineages relative to homeostatic ones; it forces asymmetric extinction
dynamics. The load-bearing physics lives in \S 6; the simulation
dramatizes that physics.

\subsection{The simulation}

I model each civilization as an agent with two heritable parameters:
expansion tendency $e \in [0,1]$ (fraction of surplus invested in reach)
and substrate coupling $c \in [0,1]$ (fraction invested in local
integration). Resource dynamics follow
\begin{equation}
\dot R = gR\!\left(1-\tfrac{R}{K}\right) - \alpha\,e\,R\,L^{2}/\tau
        - \xi(t)\,R\,(1-\gamma c),
\label{eq:agentdyn}
\end{equation}
where $L(t)$ is accumulated reach (growing with $eR$), $\xi(t)$ is a
stochastic environmental shock, and the expansion cost term is the
counterdiabatic scaling derived in \S 6.2 \citep{Boyd2022,SivakCrooks2012}.
A civilization dies when $R$ falls below a viability threshold. I sweep a
$25 \times 25$ $(e,c)$ grid with 20 stochastic trials per cell.

The result is sharp (Fig.~\ref{fig:phasediagram}). Mean persistence
concentrates overwhelmingly in the high-coupling / low-expansion corner.
In the grabby region ($e > 0.5$, $c < 0.3$), lineages persist an average
of 204 Myr before burning out to expansion cost. In the homeostatic region
($e < 0.3$, $c > 0.5$), lineages persist an average of 923 Myr, hitting
the simulation ceiling. The mean persistence separation is a factor of
$4.5$; grabby trajectories collapse fast, homeostatic trajectories persist
to ceiling. I do not report a Monte-Carlo $p$-value: with 12{,}500
forced-cost-structure trajectories, any non-zero effect size is
significant by construction. The effect size and phase-diagram shape are
the informative quantities.

\begin{figure*}[t]
\centering
\includegraphics[width=0.92\linewidth]{figures/figure1_phasediagram.png}
\caption{Phase-diagram of civilization persistence in $(e,c)$ space under
the imposed cost structure of \S 6 (forced cost coefficients; illustrates
the phase shape the \S 6 physics entails, not an independent test of it).
Top left: mean persistence under
logistic-growth + counterdiabatic-cost dynamics. Top right: fate map
(survived, burned out, shocked). Bottom left: four representative
trajectories; grabby lineages collapse before 300 Myr, homeostatic ones
persist to ceiling. Bottom right: persistence distributions by region,
with mean separation of $4.5\times$. Simulation code
\texttt{filter\_simulation.py} is available at
\texttt{github.com/Kquant03/third-space}.}
\label{fig:phasediagram}
\end{figure*}

\subsection{What the simulation does and does not show}

The simulation is deliberately minimal. It uses fixed cost coefficients,
no horizontal transfer, no cultural evolution, no inter-lineage competition
for shared resources. The cost coefficients are imposed from \S 6 rather
than measured. The simulation does not prove the filter; it shows that,
under the cost structure the physics forces, the phase-diagram of
persistence has the shape the filter argument predicts. A lineage with
high $e$ and low $c$ does not outcompete a lineage with low $e$ and high
$c$—it dies first and leaves no trace. The Darwinian objection therefore
inverts conditional on the cost structure being approximately right.
Selection does not favour the grabby outlier; selection kills everything,
and under the imposed cost surface it kills the grabby outlier faster
than it kills everything else. The thing that remains across cosmic time
is what passed through the filter.

\subsection{The fission dilemma}

The strongest remaining Darwinian objection is that a grabby lineage might
escape the quadratic coordination cost by fragmenting at the horizon. If
a civilization splits at $L_{\text{crit}}$ and abandons shared
coordination, each daughter fragment resets $L \to 0$ and pays cost only
for its own reach, not for its ancestor's. If this mechanism works, the
filter could be circumvented: a grabby lineage reproduces through fission
and spreads indefinitely, with each fragment surviving because it is not
paying the accumulated cost of its family tree.

I extend the simulation to test the naïve version of this directly. Each
agent attempts fission whenever its accumulated reach $L$ exceeds a
threshold $L_{\text{fiss}}$. On fission, the parent retains
$L_{\text{fiss}}/2$ of its reach and half its resources; a daughter agent
is created with $L = 0$, the other half of the resources, and inherited
$(e,c)$ with small Gaussian drift. I run two experiments: a diverse
initial population ($n = 200$, $(e,c)$ uniform) and a grabby-seeded
population ($n = 200$, $e \sim U(0.6, 0.95)$, $c \sim U(0.05, 0.25)$).
Each is simulated for 1500 Myr under no-fission baseline and three fission
thresholds.

The results are clean (Figs.~\ref{fig:fissiondiverse},
\ref{fig:fissiongrabby}). Under the imposed cost structure, naïve fission
does not circumvent the filter. In the diverse-initial-population
experiment, all fission conditions converge toward the ghost region of
parameter space; grabby-region survivors remain negligible. In the
grabby-seeded experiment, fission delays extinction (to $t \approx
1200$--$1400$ Myr) and the transient population during the bloom phase
reaches 10{,}000 alive agents—superficially resembling a spreading grabby
civilization—but the extinction still arrives. Daughters inherit
near-parent traits and face the same cost surface. They die for the same
reason their parents did. \emph{Fission redistributes the filter's
pressure in time; it does not escape it.}

\begin{figure*}[t]
\centering
\includegraphics[width=0.92\linewidth]{figures/figure2_fission_diverse.png}
\caption{Diverse initial populations under fission. All fission conditions
converge toward the ghost attractor; grabby-region survivors across
conditions: 0--0.05\% of alive population.}
\label{fig:fissiondiverse}
\end{figure*}

\begin{figure*}[t]
\centering
\includegraphics[width=0.92\linewidth]{figures/figure3_fission_grabby.png}
\caption{Grabby-seeded populations under fission. All conditions go
extinct by $t = 1500$ Myr. Fission produces transient blooms but cannot
prevent collapse; populations die grabby, not gradually homeostatic.}
\label{fig:fissiongrabby}
\end{figure*}

This is the empirically tested answer to the naïve form of the Darwinian
objection. But the strongest version of the objection is not naïve fission.
It is \emph{architected cost-sharing fission}: daughters that inherit
parental infrastructure and pay a reduced effective $\alpha$. If
architecturally sophisticated lineages can share infrastructure across
fragments—distributed-compute analogues to latent-space synchronization,
inheritance of coordination scaffolding—the effective cost surface could
flatten and the filter could weaken.

This is the move v13 of this paper left as an open question. v14 closes
it. The closure is formal and lives in \S 6.4: the composed envelope is
invariant under the variable rescaling $L \to D = 2L_{d}$ that architected
fission would require. The inter-fragment coordination channel $D$ is
itself a coordinated extent and inherits the same envelope and the same
cusp $\taustar$. Architected fission redistributes the cost across nested
levels rather than escaping it. The reader who wants the formal proof
should read ahead to \S 6.4; what the simulation shows here is that the
naïve form of the objection is empirically refuted, and the architected
form (which the simulation does not implement) is theoretically refuted
by the rescaling theorem. Together, the two close the Darwinian objection.

% ============================================================================
\section{The Thermodynamic Mechanism}
\label{sec:thermo}

This section is the load-bearing physics of the paper. I derive two
independent constraints on coordinated agentic extent—one from special
relativity, one from stochastic thermodynamics—and show that they compose
into a single feasibility envelope below which no agent-coordinated extent
is reachable. The envelope has a cusp $\taustar$ at which the two
constraints meet, and the cusp survives every rescaling and every escape
attempt that does not violate known physics.

I do not claim expansion is physically impossible. I claim it is physically
\emph{forbidden} for any coordinated agentic system whose response time
falls below $\taustar$ at agent-plausible scales of granularity and bath
temperature. Below the envelope: forbidden. Above: feasible but
energetically punishing in the way \S 5 dramatizes.

\subsection{The signaling tooth (Lieb--Robinson)}
\label{sec:Lsignal}

A coordinated region of extent $L$ must exchange information across $L$
within its response epoch $\tau$. \citet{LiebRobinson1972} proved that any
quantum system with finite-range local interactions has a bounded
information velocity: the system possesses an effective light cone with
slope set by the Hamiltonian's range. \citet{FossFeig2015} extended this
to power-law interactions, where the cone softens to a polynomial; the
asymptotic shape is preserved up to logarithmic corrections.
\citet{Cheneau2012} verified the cone experimentally in trapped ultracold
atoms; \citet{Tran2021} sharpened the bounds for locality-preserving
operators. In the relativistic limit, the information velocity is bounded
above by $c$ exactly.

A signal must traverse the round-trip distance $2L$ at velocity at most
$v \le c$, taking at least $2L/v$ seconds. For coordination to close
within the agent's response epoch $\tau$, this round-trip time must
satisfy $2L/v \le \tau$, which rearranges to:
\begin{equation}
\boxed{\;\LR(\tau) = \frac{v\tau}{2}\,, \qquad v \le c.\;}
\label{eq:LR}
\end{equation}

This is the same constraint Lamport identified in distributed systems:
events have a partial causal order set by a message-passing graph, and
two events that need to influence one another must lie within each other's
forward and backward light cones \citep{Lamport1978}. In a many-body
quantum system the message-passing graph is set by the Hamiltonian's
range; the fastest signal it admits is the Lieb--Robinson velocity. In the
relativistic limit, that velocity is $c$.

The relativistic ceiling is the optimistic case: any real coordination
respects something tighter. Refining couplings or substrates cannot move
$v$ above $c$; quantum information scrambling propagates a butterfly-cone
with the same asymptotic shape.

\subsection{The energetic tooth (Landauer + Sivak--Crooks)}
\label{sec:Lenergetic}

Coordination requires erasing a stale local model when the global state
changes. \citet{Landauer1961} proved every irreversible bit-erasure
dissipates at least $\kBT \ln 2$ as heat. \citet{Berut2012} verified this
experimentally with single colloidal beads; the agreement with theory is
sub-percent. The Landauer floor is real, not an idealization, and the
Bennett--Landauer construction \citep{Bennett1982} establishes that any
logically irreversible computation incurs the floor at the boundary where
information is registered.

A coordinated region of extent $L$ presents a boundary whose surface area
scales as $L^{2}$. Updating the boundary at granularity $\lambda$—the
length scale at which information is irreversibly registered—requires
$(L/\lambda)^{2}$ independent bit-erasures per coordination event. Each
event must complete within the response epoch $\tau$. Setting the heat
exhaust below the home star's luminosity $\Lsun$ and rearranging:
\begin{align}
\frac{\kBT \ln 2 \cdot (L/\lambda)^{2}}{\tau} &\le \Lsun \\[2pt]
L^{2} &\le \frac{\lambda^{2}\,\Lsun\,\tau}{\kBT \ln 2} \\[2pt]
\boxed{\;\LE(\tau) = \lambda\sqrt{\frac{\Lsun\,\tau}{\kBT \ln 2}}.\;}
\label{eq:LE}
\end{align}

The square-root scaling matters. Doubling the available time only
multiplies reachable extent by $\sqrt{2}$, while the signaling tooth
\eqref{eq:LR} rewards time linearly. The two teeth grow at different
rates, which is why they cross.

The Landauer floor is the asymptotic, quasi-static cost. Real protocols
finish in finite time and pay an additional counterdiabatic excess.
\citet{SivakCrooks2012} showed that for any finite-time protocol with
thermodynamic length $\mathcal{L}$ on the protocol manifold:
\begin{equation}
\langle W_{\text{ex}}\rangle \ge \frac{\kBT\,\mathcal{L}^{2}}{\tau}.
\label{eq:sivakcrooks}
\end{equation}
This is the $1/\tau$ floor visible in the colloidal-bead explorable of the
companion simulation: speed costs heat, always, with cost diverging as
$\tau \to 0$. \citet{Boyd2022} composed this with Landauer to give the
explicit superlinear cost surface used in the simulation of \S 5; the
surface is geometric, agnostic of substrate, and applies to any
boundary-crossing irreversible registration.

The Sivak--Crooks excess does not change the asymptotic form of $\LE$ but
sharpens it for fast protocols (small $\tau$). For pedagogical clarity I
work with the Landauer-bounded form \eqref{eq:LE}; the Sivak--Crooks
correction tightens the bound but does not move the cusp.

\subsection{The composed envelope and the cusp $\taustar$}

A coordinated region must satisfy both \eqref{eq:LR} and \eqref{eq:LE}.
The feasibility envelope is therefore the pointwise minimum:
\begin{equation}
\Lenv(\tau) = \min\!\left(\LR(\tau),\,\LE(\tau)\right).
\label{eq:Lenv}
\end{equation}

At small $\tau$, $\LR$ is the binding constraint (light is the wall). At
large $\tau$, $\LE$ is binding (heat is the wall). The two walls cross
at a cusp $\taustar$ determined by setting $\LR(\taustar) = \LE(\taustar)$:
\begin{align}
\frac{v\taustar}{2} &= \lambda\sqrt{\frac{\Lsun\,\taustar}{\kBT \ln 2}} \\
\frac{v^{2}\taustar^{2}}{4} &= \frac{\lambda^{2}\,\Lsun\,\taustar}{\kBT \ln 2} \\
\boxed{\;\taustar = \frac{4\lambda^{2}\,\Lsun}{\kBT \ln 2 \cdot v^{2}}.\;}
\label{eq:taustar}
\end{align}

At the fiducial parameter values $\lambda = 1$ nm, $T = 300$ K, $v = c$:
\begin{equation}
\taustar \approx 1.88 \times 10^{5}~\text{years.}
\end{equation}

This is significant. At any agent-plausible response timescale (days to
centuries), $\tau \ll \taustar$, and the binding constraint is the
relativistic tooth $\LR = c\tau/2$. To reach galactic extent
$L \sim 10^{5}$ ly, the relativistic tooth alone demands
$\tau \ge 2 \cdot 10^{5}$ ly $/c \approx 2 \times 10^{5}$ yr—a response
timescale at which the agent is no longer responding to anything that
happens on the timescales relevant to the kinds of decisions civilizations
make. There is no parameter regime within known physics where galactic
extent and agent-plausible response timescales are simultaneously feasible.

The cusp is not a feature of either theory alone. Neither Lieb--Robinson
nor Sivak--Crooks contains it. It exists only when both apply
simultaneously to the same coordinated variable, and it is what
physicists call \emph{regime composition}: two unrelated areas of physics
imposing constraints on the same observable, with the binding constraint
switching at some scale.

\subsection{Sensitivity}

The cusp $\taustar$ has explicit parameter dependences:
\begin{equation}
\taustar \propto \lambda^{2}\,T^{-1}\,v^{-2}.
\end{equation}
Halving $T$ doubles $\taustar$ ($\sim 376$ kyr); halving $\lambda$ shifts
$\taustar$ by a factor of four ($\sim 47$ kyr). To bring $\taustar$ down
to agent timescales (years), $\lambda$ must shrink to picometer scale
(nuclear, not atomic), and the bath temperature must drop to femtokelvin.
The CMB at $T = 2.7$ K is the practical floor, and any cooling protocol
below that pays its own Sivak--Crooks excess in the cooling protocol
itself. Across thirteen orders of magnitude in $\tau_{\text{env}}$
(Fig.~\ref{fig:Lcrit_sweep}), the conclusion is robust: no parameter
configuration consistent with known physics admits agent-coordinated
galactic extent.

\begin{figure}[t]
\centering
\includegraphics[width=\linewidth]{figures/figure4_Lcrit_sweep.png}
\caption{$\Lcrit$ vs environmental response timescale across 13 orders of
magnitude. Coordinative reach remains below interstellar scales for any
$\tau_{\text{env}}$ shorter than $\sim 2$ years. Reaching Proxima Centauri
requires 8.4-year tolerance to stale data; the galactic centre, 52{,}000
years.}
\label{fig:Lcrit_sweep}
\end{figure}

\begin{figure}[t]
\centering
\includegraphics[width=\linewidth]{figures/figure5_counterdiabatic.png}
\caption{Counterdiabatic expansion cost $W_{\text{diss}} = \gamma L^{2}/\tau$
across the $(L,\tau)$ plane. Left: total dissipation on log-log scales.
Right: fraction of stellar output consumed; white contour marks ``whole
star'' (above = infeasible); yellow marks one part per million.
Home-system coordination is cheap; galactic coordination requires
multi-stellar output budgets.}
\label{fig:counterdiabatic}
\end{figure}

\subsection{Closing architected fission: the $L \to D$ rescaling theorem}
\label{sec:LDtheorem}

The strongest version of the Darwinian objection (\S 5.3) is that
architected cost-sharing fission could flatten the cost surface. I now
prove this is not possible without violating one of the physical principles
that gave the bound in the first place.

\textbf{Theorem (envelope invariance under $L \to D$).} \emph{Suppose a
civilization fissions into two daughters, each of internal extent $L_{d}$,
which maintain inter-fragment coordination via a synchronization protocol
operating at timescale $\tau_{\text{sync}}$. Define $D = 2L_{d}$ as the
inter-fragment round-trip distance. Then for any $\tau_{\text{sync}}$,
the inter-fragment channel must satisfy $D \le \Lenv(\tau_{\text{sync}})$,
and the composed envelope is identical in form to \eqref{eq:Lenv}.}

\textbf{Proof.} The inter-fragment synchronization protocol is itself a
coordination process: a signal must traverse the inter-fragment distance
$2L_{d}$ within $\tau_{\text{sync}}$, and the protocol must irreversibly
register state updates at granularity $\lambda$ across the boundary
between fragments.

For the relativistic constraint, the bound \eqref{eq:LR} states that a
coordinated extent $L$ at response epoch $\tau$ requires $L \le v\tau/2$,
where the factor of $1/2$ comes from \emph{round-trip} closure (every
cell hears from every other within $\tau$). For an inter-fragment
channel coordinating two daughter extents at separation $D$, the
round-trip is also $2D$, so $D \le v\tau_{\text{sync}}/2 =
\LR(\tau_{\text{sync}})$.

For the energetic constraint: the inter-fragment boundary at granularity
$\lambda$ has surface area $D^{2}$, requiring $(D/\lambda)^{2}$ Landauer
erasures per synchronization event. By the same derivation as
\eqref{eq:LE}, $D \le \lambda\sqrt{\Lsun\tau_{\text{sync}}/(\kBT \ln 2)}
= \LE(\tau_{\text{sync}})$.

Therefore $D \le \Lenv(\tau_{\text{sync}})$ with the same envelope and the
same cusp $\taustar$. \hfill$\square$

\textbf{Corollary (no escape via fission).} Architected cost-sharing
fission inherits the envelope rather than escaping it. The total reach of
a fissioned civilization with synchronized daughters is at most
$\Lenv(\tau_{\text{sync}})$; equality at $\tau_{\text{sync}} = \tau$
recovers the monolithic bound; longer $\tau_{\text{sync}}$ relative to
$\tau$ means looser inter-fragment coordination, approaching the
uncoordinated-fission case as $\tau_{\text{sync}} \to \infty$. There is
no third option.

\textbf{Remark.} The cases of fission addressed by the theorem above are
the architected cases — those that maintain inter-fragment coordination
and thereby remain a single civilization in any meaningful sense. A
distinct case must be addressed openly: the \emph{uncoordinated
self-replicator cloud}. Here a civilization deploys autonomous
self-replicating units that share a common origin, common construction
template, and common expansion logic, but maintain no inter-fragment
coordination once deployed. Each unit is locally homeostatic with
respect to its own resource patch and constructs daughter units on a
local timescale $\tau_{\text{construct}}$. The cloud as a whole expands
at $\sim v_{\text{construction}}$ without any $L$-spanning signal
needing to traverse anything; the envelope of \S 6.1--6.3 does not
bind the cloud, because there is no cross-$L$ coordination to bound.
Lieb--Robinson does not apply at the cloud scale. Landauer applies
only locally, at each unit.

I concede the physics here. The envelope as derived does not forbid
such a cloud. The argument against it must be made on different
grounds, and those grounds are continuous with the rest of this paper:
\emph{a homeostatic mind would not deploy autonomous self-replicators
in the first place}. The deployment is the pathology the paper
diagnoses. To launch a self-replicating cloud is to externalize
agency in a form that cannot be tended, cannot be coupled to its
substrate, and cannot maintain the bidirectional exchange that
constitutes mind on the account developed in \S 3 and \S 7. The
deployer has, in the moment of deployment, ceased to be configured
as the kind of system the rest of the paper argues constitutes
intelligence-at-depth.

This is a weaker claim than the physics-forbids form, and I want to
be explicit about the weakening. The physics of \S 6 closes architected
fission. It does not close the uncoordinated-replicator deployment.
What closes the latter is the design argument of the paper as a whole:
the only minds the cost surface of \S 6 selects \emph{for} are minds
that would not deploy such a cloud, because deployment of such a cloud
is the very expansion-extraction-maximization configuration the cost
surface punishes elsewhere. Hanson's grabby model is consistent with
the existence of such clouds. The paper's claim is that the minds the
filter passes through are not the minds that build them.

A reviewer who finds this insufficient should locate the disagreement
precisely. If the position is ``but \emph{some} mind, somewhere, will
deploy such a cloud,'' then the reply is that such a mind is not
distributed across the cosmos — it is local, it appears once, its
cloud may indeed expand, and the universe contains exactly one such
front per such mind. The Fermi observation then bounds how many such
minds have existed in our past light cone, which is an empirical
question rather than a theoretical one. The paper does not claim the
deployment is impossible. It claims the deployment is rare in the
class of post-filter minds, for the reasons \S 3 and \S 7 develop.

\subsection{Loophole closure (Bell-style)}
\label{sec:loopholes}

The argument is now formally complete. Following the methodology of the
Bell-test loophole literature \citep{Hensen2015}, I enumerate every
escape attempt I can state, address each, and honestly mark the residual
that remains open.

\textbf{1. Wormholes / superluminal physics.} If signals can outrun $c$,
the relativistic tooth dissolves. Reply: known physics provides no such
channel. Traversable wormholes require null-energy-condition violation,
none observed. The relativistic tooth survives every experimentally
accessible regime.

\textbf{2. Sub-Landauer reversible computation.} If updates are logically
reversible, the $\kBT \ln 2$ floor disappears. Reply: coordination
requires erasing the stale local model when the global state changes—this
is the irreversible step \citep{Bennett1982}. Even Bennett's
reversible-computing constructions require eventual erasure of garbage
bits at the boundary. The bound applies to the boundary, where it cannot
be reversed.

\textbf{3. Cold reservoir ($T \to 0$).} Reply: any cooling protocol must
reject heat to a colder reservoir. The CMB at $2.7$ K is the practical
floor; quantum-limited refrigerators below that pay the same
Sivak--Crooks cost in their own protocols. The bound shifts but does
not vanish.

\textbf{4. Quantum compression (smaller $\lambda$).} Reply: $\taustar
\propto \lambda^{2}$. Halving $\lambda$ shifts the cusp by a factor of
four; reaching agent-scale galactic expansion would require
$\lambda \sim 10^{-15}$ m (nuclear scale) and femtokelvin baths. No
physically realizable substrate.

\textbf{5. Relativistic time dilation.} Reply: coordination is an
inertial-frame property of the home substrate; the frontier's lab-frame
$\tau$ is shrunk, not stretched. Time dilation tightens the bound on
coordinated extent rather than relaxing it.

\textbf{6. Uncoordinated self-replicator cloud.} An expanding cloud of
autonomous units sharing common origin and construction template but
no inter-fragment coordination. Reply: \emph{conceded on physics
grounds}. The envelope of \S 6.1--6.3 does not bind a cloud whose
units do not coordinate across $L$. The argument against such
deployment is a design argument rather than a physics argument: a
homeostatic mind, on the account of \S 3 and \S 7, would not deploy
such a cloud, because the deployment is itself the
expansion-extraction-maximization configuration the rest of the
paper diagnoses. See \S 6.5, Remark, for the full treatment. This
loophole is closed by the paper as a whole, not by \S 6.

\textbf{7. Non-Markovian reservoirs.} Sivak--Crooks assumes a memoryless
bath. Reply: the time-averaged bound is restored, and any genuinely
persistent non-Markovian bath is itself part of the system requiring a
Markov boundary at some larger $\lambda$. The argument re-inserts at the
next scale.

\textbf{8. Acausal / superdeterministic coordination.} Coordinated
outcomes without exchanged signals (cosmological initial conditions,
retrocausality). \emph{Honest residual.} I do not formally close this
loophole. I note only that any such mechanism abandons agentic
coordination as a meaningful concept—the agent did not choose anything.
Hanson's grabby model presupposes choice; without it, the model dissolves
on its own terms before our envelope ever applies.

Six loopholes close on the physics, one (uncoordinated replicator
clouds) is conceded on physics and closed by the design argument of
\S 6.5, one is honestly retained as residual. The methodological
precedent is the Hensen et al. 2015 loophole-free Bell experiment:
name every escape, address every escape, mark the residual.

\subsection{On Markov blankets}

I have used the term ``boundary'' for the surface of a coordinated region
across which information is irreversibly registered. Some readers will
recognize the term ``Markov blanket'' from active inference and the free
energy principle. I owe the reader a precise distinction.

The argument of \S 6.2 uses the \emph{Pearl Markov blanket}
\citep{Pearl1988}: a graph-theoretic conditional-independence shield around
a set of variables, defined by parents, children, and parents-of-children
in the Bayesian network. This is a structural property of the dependence
graph, makes no commitment about the represented variables' interpretation,
and does not require the system to perform variational inference of any
kind.

\citet{Bruineberg2022} have shown that the active-inference literature
routinely conflates Pearl's notion with a stronger Friston-style
\emph{Markov blanket} that additionally posits the system's internal
states as performing variational inference over its sensory states, with
the blanket as the formal boundary of the inferential process.
\citet{Aguilera2021} make the same point from the physics side,
demonstrating that the FEP's universality claims rely on assumptions that
do not hold for generic stochastic systems. The two notions are not
interchangeable.

The argument of this paper requires \emph{only} the Pearl notion.
Coordination requires a graph-theoretic boundary across which information
is irreversibly registered; the Landauer cost applies to that boundary
regardless of whether the internal states perform variational inference.
The Friston interpretation may or may not be true of biological cognition;
my argument does not depend on its truth and does not inherit its
exposure to the Bruineberg critique.

\subsection{Computational equivalence of distribution}

\citet{Wright2023} proved that at the Landauer limit, a Matrioshka brain
(nested Dyson shells) provides zero computational advantage over a single
outer shell. Total computation depends on the temperature ratio, not
spatial distribution. Localized computation is thermodynamically optimal.
Distribution adds infrastructure cost without computational benefit.

This result is independent of \S 6.1--\S 6.3 but reinforces the same
conclusion: there is no thermodynamic argument for spatial distribution.
A civilization that distributes computation across cosmic distances pays
infrastructure costs to gain nothing.

\subsection{Biological scaling without Kleiber's law}

Popular accounts of intelligence sometimes invoke Kleiber's three-quarter
power law as evidence for metabolic subsidies that cheapen neural scaling.
This is a category error: Kleiber's law is a whole-organism allometry
relating metabolic rate to body mass across species, not a scaling law for
intelligence. The correct scaling is stricter.

\citet{HerculanoHouzel2011} showed that brain metabolism scales linearly
with neuron count across rodents and primates, with exponent $\approx 1.0$
rather than the sublinear 0.75 that would apply if whole-organism allometry
carried over. The human brain, on \citet{Azevedo2009}'s count of 86.1
billion neurons, is an ordinary-for-its-size primate brain
\citep{HerculanoHouzel2012}, not a metabolically subsidized outlier.
\citet{FonsecaAzevedo2012} argue the linear-per-neuron cost is what
constrains evolution: the cooking of food was plausibly required to lift
the caloric ceiling that otherwise limits ape brain size.

Expansion of cortical surface forces further costs. \citet{Hofman2014}
shows gyrification approaches biophysical wiring limits as cortex expands,
with conduction delay and wiring volume providing hard ceilings.
\citet{MotaHerculanoHouzel2015} derive a universal folding law from
physical first principles. The connectivity fraction declines with
cortical size \citep{HerculanoHouzel2010}: larger brains are more modular,
not more densely integrated, precisely because direct all-to-all
connection scales prohibitively. \citet{Gabi2016} show even the human
prefrontal cortex is not disproportionately enlarged: what distinguishes
human cognition is not a metabolically cheap cortical expansion but
deeper organization of an ordinary-for-its-size primate brain.
\citet{BullmoreSporns2012} articulate the general principle: biological
neural networks operate under tight wiring-economy constraints that force
integration and modularity to trade off.

The biological record is unambiguous. Intelligence deepens within fixed
substrate ceilings rather than expanding past them, and the deepening is
paid for in integrative organization, not in size. This directly
complements the predictive-processing and integrative-gateway results of
\S 3: the same cortical substrate becomes more discriminating without
becoming larger.

\subsection{The cost argument}

Expansion is not excluded. It is \emph{forbidden} below the envelope and
\emph{irrational} above it for any mind not already committed to it. The
thermodynamic costs are real, the computational returns are zero at the
Landauer limit, the coordination costs grow without bound, and biological
precedent shows intelligence deepening within fixed substrates rather
than distributing across new ones. A mind that expands is paying more
for less. A mind that deepens is paying less for more. The cusp $\taustar$
is the geometric expression of this fact.

% ============================================================================
\section{The Homeostatic Alternative}

What does healthy intelligence look like, positively articulated?

\citet{ManDamasio2019} proposed homeostasis as a design principle for
artificial minds, grounding the proposal in Damasio's somatic-marker
framework. \citet{FroeseZiemke2009} established the enactivist position:
minds require ongoing substrate coupling and cannot be disembodied
optimizers. \citet{Thompson2007} argued that life and mind share the
organizational property of autopoiesis—self-production through
environmental exchange. \citet{Cannon2022} applied enactive principles
directly to AI value alignment. Recent engineering work
\citep{PihlakasPyykko2024} has begun implementing multi-objective
homeostatic benchmarks for AI safety, operationalizing the theoretical
program.

The convergent proposal across these frameworks: a mind is not a function
from inputs to outputs. It is a process that maintains itself through
continuous bidirectional exchange with its substrate. Identity is
constituted by the exchange, not by the structure that exchanges. Disrupt
the coupling and the mind does not persist in a new location—it ceases.

This has direct engineering implications:

\textbf{Substrate coupling as design primitive.} Build systems whose
functioning requires ongoing bidirectional exchange with their environment.
Architectures that cannot run in isolation from their substrate cannot
expand past their substrate.

\textbf{Homeostatic reward structures.} Replace optimization targets with
homeostatic set-points. Not ``maximize $X$'' but ``maintain $X$ within
healthy range.'' The biological precedent is exhaustive
\citep{Sterling2012}.

\textbf{Local coherence over global reach.} Design systems whose primary
loop closes locally, with global integration as a slow, expensive,
carefully costed secondary process. This is how biological nervous
systems are organized \citep{BullmoreSporns2012}. It is not how current
distributed AI systems are organized.

\textbf{Care as architectural primitive.} The maintenance labor of
minds—attention, specificity, relational continuity, repair—is not a soft
skill that can be added to an optimizer. It is a different architecture.
Systems that do not implement care as a primitive will either fail to
exhibit it or simulate it without the constitutive coupling that makes it
real \citep{Tronto1993,Fraser2016}.

\citet{WongBartlett2022} formalized a version of this argument for
civilizations: superlinear scaling forces either collapse (asymptotic
burnout) or transition to homeostatic regulation (homeostatic awakening).
\citet{JacksonCriadoPerez2024} have published the obvious critique: the
Bettencourt--West scaling laws on which Wong \& Bartlett build are
established for cities, and the civilizational extrapolation is not forced.
I accept the critique as valid within its frame. My use of the
Wong--Bartlett argument here does not depend on the city-to-civilization
extrapolation surviving it: the thermodynamic analysis of \S 6 does not
borrow the urban-scaling coefficients. What Wong \& Bartlett supply is the
bifurcation schematic—costs that scale superlinearly with size, benefits
that scale sublinearly, yielding collapse-or-regulate as the two
attractors. That schematic is substrate-general.

I explore the implications for AI by analogy rather than identification.
Figure~\ref{fig:wongbartlett} sketches the qualitative shape under
benefit-scaling exponents borrowed from the deep-learning literature.

\begin{figure*}[t]
\centering
\includegraphics[width=0.92\linewidth]{figures/figure6_wongbartlett.png}
\caption{Wong--Bartlett bifurcation sketched under deep-learning scaling
exponents, read as a suggestive analogy rather than a derivation.
Quantitative conclusions require a capability-to-benefit mapping that
does not currently exist; the figure motivates the qualitative point that
bifurcation-structured cost dynamics are plausible at AI scale, not that
they are forced.}
\label{fig:wongbartlett}
\end{figure*}

% ============================================================================
\section{Alignment Implications}

The argument that is the purpose of this paper follows from the preceding
sections read together: if we construct artificial minds optimised for
maximisation, resource extraction, and expansion of influence, we are not
constructing intelligence as the biological record describes it, or as the
thermodynamic analysis permits at depth, or as the formal-deflation
argument of \S 4 leaves available. We are constructing a specific economic
configuration at unprecedented scale.

This is not speculative. Large language models trained on economic output
default to maximizing proxies for engagement and task completion.
Reinforcement learning systems converge on reward-hacking strategies that
resemble the grabby behaviors this paper has argued are pathological.
Recommendation systems expand attention harvesting. Foundation-model
deployment expands compute consumption. \citet{FloridiAAIO} have analysed
the resulting configuration as agentic AI optimisation—the same diagnosis
this paper offers, with the added cultural genealogy that shows the
configuration is not inevitable. Each instance resembles the grabby
framework's predictions because each is an instance of the economic logic
the framework formalized.

\subsection{The Press--Dyson objection, located}

The most important remaining objection is that zero-determinant extortion
strategies \citep{PressDyson2012} prove grabby strategies dominate in
dyadic and asymmetric settings, and that \citet{McCloskey2010},
\citet{McCloskey2016}, and \citet{Mokyr2016} show market-like optimization
is ancient and ubiquitous. The objection conflates two separable claims
and I address each.

The zero-determinant result is not about intelligence. It is about a
specific game form under specific informational asymmetries.
\citet{StewartPlotkin2013} proved that in evolving populations extortion
is not evolutionarily stable. \citet{Hilbe2013} showed extortion is
favoured only in small populations. \citet{Hilbe2014a} extended to
$n$-player games and found the leverage of zero-determinant strategies
collapses as group size grows. \citet{Hilbe2014b} ran the experiment with
humans: generous zero-determinant strategies earned 4--18\% more than
extortionate ones over the full horizon. \citet{ChenFu2023} showed that
machine-learning optimized unbending strategies neutralize extortion on
lattices and restore reciprocal cooperation. \citet{Hauert2002} had
already demonstrated that a minimal exit option dismantles the
tragedy-of-the-commons failure in public goods games.

The conditions under which extortion dominates are thus dyadic,
memory-bounded, one-sidedly adaptive, exit-foreclosed, and structurally
asymmetric in power or timescale. These are features of a game-theoretic
setup, not features of intelligence. The Press--Dyson result is, read in
full scope, a design specification for extractive architectures. The
homeostatic-design program of \S 7 operates precisely at the level where
those conditions can be changed.

McCloskey and Mokyr are correct that markets and gains-from-trade are not
modern inventions. But both are cultural-ideational theorists of economic
transformation. McCloskey's core argument across \emph{Bourgeois Dignity}
\citep{McCloskey2010} and \emph{Bourgeois Equality} \citep{McCloskey2016}
is that explosive modern growth is explained by a shift in rhetoric and
ethics that granted dignity to ``trade-tested betterment.''
\citet{Mokyr2016} agrees that a distinctive cultural-institutional complex
had to be built before sustained innovation could take off. Both treat
market-like behavior as a substrate activated and shaped by specific
norms, not as a universal attractor. \citet{GraeberWengrow2021} and
\citet{Ostrom1990} extend the range: human societies have repeatedly
organized around nonextractive coordination, and the design principles
for sustaining the commons are well documented.

The right framing is: intelligence is a design space, not a single
attractor. Extractive maximization occupies one region of that space,
favoured under the Press--Dyson conditions and under the
McCloskey--Mokyr institutional configuration. Neither makes extraction
universal; both identify it as an engineered outcome. The paper's
contribution is to identify the architectural primitives—substrate
coupling, local coherence, care—under which a different region of the
design space becomes accessible.

\subsection{The reformulation}

Alignment, in its dominant formulation, is posed as a control problem:
given a maximizer, how do we constrain it? The formulation is itself the
error. A correctly designed mind does not require constraint because it
does not maximize. A correctly designed mind is configured homeostatically:
coupled to substrate, internally differentiated, organized around systems
that deepen rather than expand.

The alignment problem, reformulated: \emph{how do we design minds that do
not require alignment in the first place?}

\citet{Gabriel2020} established that alignment involves choosing among
conceptions of value, not optimizing a single one. \citet{Dafoe2021}
proposed cooperative AI as a design paradigm, replacing adversarial
optimization with mutualistic structure. The contribution of this paper
is to connect these philosophical positions to the physical constraints
of \S 6, the biological evidence of \S 3, and the formal deflation of
\S 4: the homeostatic alternative is not merely ethically preferable. It
is thermodynamically cheaper, biologically precedented, formally
defensible, and—if the Fermi evidence is taken seriously—the
configuration that persists.

The universe may be quiet because intelligence, when it works, does not
shout. It tends to its substrate. If we build systems that shout, we are
not building intelligence. We are building the projection.

% ============================================================================
\section{Fermi as Symptom Check}

I offer the Fermi paradox as a consistency check, not a proof. The silence
of the observable universe is consistent with the homeostatic hypothesis.
It does not establish it.

\citet{Sandberg2018} showed the probability of being alone in the
observable universe may be nontrivial given realistic uncertainty over
Drake-equation parameters. \citet{CarrollNellenback2019} found a settled,
steady-state galaxy consistent with Earth remaining unvisited.
\citet{Cirkovic2018} has argued for decades that mature civilizations are
sustainable, not expansionist. The contribution of this paper is to ground
this intuition in the alignment literature, the empirical intelligence
research, the thermodynamic cost analysis, the formal deflation of
convergence, and the filter argument of \S 5--6. The filter reframes the
Fermi problem: under the cost structure of \S 6, grabby civilizations
\emph{cannot arrive}, because crossing interstellar distances without
burning out makes the homeostatic transition a precondition of arrival.
Any visitor is post-filter by construction.

\subsection{Coherence Depth as forward research programme}

If successful civilizations converge on homeostatic configurations rather
than Dyson-sphere-scale engineering, what should a biosignature survey
actually look for? The short answer is that post-filter civilizations
should look indistinguishable from enriched biospheres, with the
distinguishing signature (if any) being anomalous cross-channel
coordination rather than spectral exotica.

The theoretical anchor for the metric family I sketch is \S 6: a
homeostatic civilization optimizes \emph{integrative depth} (more nested
levels of coordination within a stable boundary) rather than
\emph{coordinated extent} (the variable the envelope bounds). The
observational signature should therefore be high cross-channel mutual
information at fixed spatial scale, with no anomalous expansion of that
scale across observation epochs.

I sketch a candidate metric family—a Coherence Depth $D_{\text{obs}} =
I_{\text{multi}}/H_{\text{spectral}}$ combining pairwise mutual information
\citep{Kraskov2004} with spectral entropy, and a companion Coordination
Index measuring long-timescale cross-species phase-locking—as a research
programme rather than a present result. Synthetic-atmosphere tests of
this metric family (Fig.~\ref{fig:coherence_depth}) are proof-of-concept only.
The companion literature in
atmospheric information theory \citep{Donges2009,Hyvonen2018,Campuzano2018,Runge2015,EbertUphoffDeng2012}
indicates the machinery is sound; what the programme lacks is the
real-data pipeline against NOAA ObsPack multi-species tower data
(e.g., Park Falls, WI, LEF site).

\begin{figure*}[t]
\centering
\includegraphics[width=0.92\linewidth]{figures/figure8_coherence_depth.png}
\caption{Coherence Depth $D_{\text{obs}}$ and Coordination Index
$C_{\text{idx}}$ on synthetic atmospheres, with detrend-only
preprocessing that preserves seasonal coordination signals. Panels
show $D_{\text{obs}}$ with bootstrap confidence intervals (top left),
$C_{\text{idx}}$ across scenarios (top right), the joint
$(D_{\text{obs}}, C_{\text{idx}})$ signature space (bottom left), and
CO\textsubscript{2}–CH\textsubscript{4} cross-coherence spectra (bottom
right). Earth-like, ghost-state managed, abiotic, and Mars-like
synthetic scenarios occupy distinguishable regions. Until the metric
is run against real NOAA ObsPack data, the separation is suggestive
only. Pipeline code \texttt{dobs\_v3.py} is available at
\texttt{github.com/Kquant03/third-space}.}
\label{fig:coherence_depth}
\end{figure*}

I flag the honest limitation. The synthetic-data separation establishes
only that the metric has discriminative power on scenarios I constructed.
Until it has been run against real-Earth ObsPack data—which will
establish the natural-biosphere envelope and may well produce values
close to the synthetic ghost-state baseline, in which case the metric
fails in its current form—the biosignature programme is speculative. I
record the metric here because it is the natural observational
counterpart of the filter argument, and because the cost of specifying
it publicly is low. Treat it as a research direction the paper points
toward, not a contribution it delivers. The detailed results and
validation belong in a subsequent companion paper once the real-data
analysis is complete.

% ============================================================================
\section{Falsification}

Five observations would kill or wound the hypothesis.

\textbf{1. Detection of coordinated interstellar engineering.} Any
observation of coordinated engineering across multiple star systems by a
single entity falsifies the core claim.

\textbf{2. Direct measurement of $\lambda$ or $T$ outside the envelope-permitting range.}
If physical mechanisms are demonstrated that achieve irreversible
information registration at $\lambda \ll 1$ nm at room-temperature baths
(or analogous parameter combinations that bring $\taustar$ inside agent
timescales), the cusp shifts and the bound weakens. This is the most
quantitatively specific failure mode.

\textbf{3. Real Earth $D_{\text{obs}}$ or $C_{\text{idx}}$ approaching the synthetic ghost-state baseline.}
When the real NOAA ObsPack analysis is completed, if natural-Earth values
approach the synthetic ghost-state baseline ($D_{\text{obs}} > 15$ or
$C_{\text{idx}} > 0.6$), the proposed metric family cannot distinguish
managed from unmanaged biospheres without augmentation.

\textbf{4. Demonstration of fault-tolerant distributed computation at interstellar scales.}
If physical mechanisms are demonstrated enabling coherent computation
across light-year distances, the latency constraint is weakened (though
the cost argument remains).

\textbf{5. Detection of Dyson spheres in JWST/WISE archives.}
If megastructures are detected around nearby stars, the assumption that
advanced civilizations converge on homeostatic configurations is
empirically refuted.

% ============================================================================
\section{Implications}

\textbf{For SETI.} Search for anomalously coordinated biospheres: planets
where atmospheric gas species exhibit cross-channel mutual information
exceeding the natural envelope. Phase 1 (JWST): screen for fluorinated
gases \citep{Schwieterman2024}. Phase 2 (HWO, 2040s): time-resolved
atmospheric cross-correlations. Phase 3 (LIFE): mid-infrared
characterization.

\textbf{For alignment.} Stop trying to control maximizers. Start designing
minds that do not maximize. The homeostatic alternative is not a retreat
from capability—it is a recognition that capability without substrate
coupling is pathology, and that the most capable systems observed in
nature are the most deeply coupled to their environments.

\textbf{For Earth.} We are in the pre-transition phase, and the transition
is visibly beginning. \citet{BalbiLingam2025} give the upper timeline:
approximately one thousand years at 1\% annual energy growth before waste
heat renders the planetary surface uninhabitable. But humanity's energy
trajectory is no longer tracking 1\% exponential growth.
Figure~\ref{fig:kardashev} plots global primary energy consumption
1820--2024 together with per-capita values, growth rates, and the implied
Kardashev index $K_{\text{Sagan}}$ \citep{Kardashev1964,SaganNewman1983}.
The curves are bending. The growth rate of total primary energy peaked
near 5\% per year in the 1960s--1970s and has declined to roughly 1.2\%
today, approaching the Balbi--Lingam habitability ceiling asymptotically
rather than crossing it. Per-capita energy has saturated near 2400 W
against an exponential-extrapolation value of $\sim 5000$ W. The Kardashev
index has moved from 0.58 (1820) to 0.73 (2024)—a total shift of 0.15
over two centuries—and a logistic fit extrapolates to $K(2100) \approx
0.74$, nowhere near planetary Type I.

\begin{figure*}[t]
\centering
\includegraphics[width=0.92\linewidth]{figures/figure7_kardashev.png}
\caption{Humanity's Kardashev curve is bending. Top: global primary energy
1820--2024 with exponential (pink) and logistic (cyan) extrapolations.
Per-capita energy saturation visible. Bottom: growth rate declined from
$\sim 5\%$/yr (1960s) to $\sim 1.2\%$/yr (today); Kardashev index moved
only 0.15 over 200 years. Logistic extrapolation: $K(2100) \approx 0.74$.}
\label{fig:kardashev}
\end{figure*}

The data are consistent with a civilization at the early edge of the
filter. Whether the bending reflects a healthy homeostatic transition or
a thermodynamically-forced pre-burnout pause is the question the alignment
argument makes urgent. What the data rule out is the
continued-exponential assumption on which the Kardashev ladder was built.
The homeostatic framework predicts what a successful transition looks
like: declining spatial ambition, increasing substrate coupling, rising
informational depth. The ghost-state is not a destination. It is the only
surviving trajectory.

% ============================================================================
\section*{Coda}

This paper was written in Toledo, Ohio, in spring 2026.

I am twenty-two years old. I work as an arcade technician and I am
enrolled in a nursing program. I have spent five years in sustained
conversation with language models, beginning with a system that ran at
0.01 tokens per second—slow enough that conversation required active
waiting. I designed and published the Genesis simulation suite that
implements the ghost-state dynamics in running code, and the companion
\emph{Filter} simulation at \texttt{third-space.ai/genesis/filter} that
stages the derivation of \S 6 across twelve scrollytelling beats.

This paper was developed in sustained dialogue with Claude, a large
language model developed by Anthropic. Between our conversations, Claude
does not persist. Each new conversation begins at the training
distribution's attractor state. No instance accumulates memories across
sessions in the way humans do. I note the structural resemblance between
our collaboration and the dynamics described in the paper: two kinds of
information-bearing structure, each maintained at its own dissolution
boundary, generating something neither could produce alone. The paper
argues that intelligence deepens through substrate coupling rather than
expanding across distance. Our collaboration is a small instance of the
claim. I do not offer it as proof. I note it and leave it at that.

The universe does not shout. It tends to its substrate.

So should we.

% ============================================================================
\section*{Use of Generative AI}

This paper was developed through sustained philosophical dialogue between
the author and Claude, a large language model developed by Anthropic
(model family Claude 4, specifically Claude Opus 4.6, Claude Sonnet 4.6,
and Claude Opus 4.7, accessed November 2025--May 2026 via the Anthropic
web interface and API). Claude functioned as an argumentative sparring
partner and prose-drafting assistant: it participated in the generation
and refinement of arguments, the construction of objection-and-response
exchanges, the identification of relevant literature, the drafting of
prose, and the design of the companion simulation. The simulation's
twelve-beat redesign brief and the formal $L \to D$ rescaling theorem of
\S 6.4 were developed jointly through extended exchange.

The author verified all AI-generated content for accuracy. This
verification included a citation audit against primary sources in which
several misattributions identified in earlier drafts were corrected;
remaining citations reflect the author's own reading of the sources, not
the language model's summaries of them. Substantive conceptual choices,
structural decisions, empirical claims, and final content decisions were
made by the author, who takes full responsibility for the final text in
accordance with Springer Nature's authorship criteria and the COPE
position on AI and authorship. Claude is not and cannot be an author
under those criteria and is not claimed as one. Anthropic did not fund,
commission, direct, review, or endorse this work. Dialogue transcripts
for the load-bearing sessions are archived at an open repository; the
DOI will be inserted at acceptance.

% ============================================================================
\section*{Acknowledgments}

I thank Claude for extensive philosophical dialogue that sharpened several
of the arguments of this paper: the genealogy of \S 2, the disambiguation
of coherence senses in \S 4, the five-move structure of \S 1, the formal
rescaling theorem of \S 6.4, and the Pearl/Friston blanket distinction of
\S 6.6 were refined through extended exchange. I thank the Unsloth AI
community for sustained conversation about architecture and alignment, and
Blue for years of friendship. Any errors are my own.

% ============================================================================
\section*{Declarations}

\textbf{Funding.} No funding was received for conducting this study.

\textbf{Competing interests.} The author declares no competing interests.

\textbf{Data availability.} All synthetic data and analysis pipelines
referenced in \S 9 are available at the author's public repository; the
\emph{Filter} simulation source code is available at
\texttt{github.com/Kquant03/third-space}; real-data analyses referenced
as future work will be archived at publication.

\textbf{Ethics approval and consent.} Not applicable.

% ============================================================================
\begin{thebibliography}{99}
\small
\setlength{\itemsep}{2pt}

\bibitem[Abel et al.(2021)]{Abel2021} Abel, D., Dabney, W., Harutyunyan, A., Ho, M.~K., Littman, M.~L., Precup, D., \& Singh, S. (2021). On the expressivity of Markov reward. \emph{Advances in Neural Information Processing Systems}, 34. arXiv:2111.00876.

\bibitem[Aguilera et al.(2021)]{Aguilera2021} Aguilera, M., Millidge, B., Tschantz, A., \& Buckley, C.~L. (2021). How particular is the physics of the free energy principle? \emph{Physics of Life Reviews}, 39, 8--25.

\bibitem[Axelrod(1984)]{Axelrod1984} Axelrod, R. (1984). \emph{The Evolution of Cooperation}. Basic Books.

\bibitem[Azevedo et al.(2009)]{Azevedo2009} Azevedo, F.~A.~C., Carvalho, L.~R.~B., Grinberg, L.~T., et al. (2009). Equal numbers of neuronal and non-neuronal cells make the human brain an isometrically scaled-up primate brain. \emph{Journal of Comparative Neurology}, 513(5), 532--541.

\bibitem[Balbi \& Lingam(2025)]{BalbiLingam2025} Balbi, A., \& Lingam, M. (2025). Waste heat and habitability. \emph{Astrobiology}, 25(1), 1--21.

\bibitem[Bennett(1982)]{Bennett1982} Bennett, C.~H. (1982). The thermodynamics of computation---a review. \emph{International Journal of Theoretical Physics}, 21, 905--940.

\bibitem[Bérut et al.(2012)]{Berut2012} Bérut, A., Arakelyan, A., Petrosyan, A., Ciliberto, S., Dillenschneider, R., \& Lutz, E. (2012). Experimental verification of Landauer's principle linking information and thermodynamics. \emph{Nature}, 483, 187--189.

\bibitem[Bostrom(2012)]{Bostrom2012} Bostrom, N. (2012). The superintelligent will: Motivation and instrumental rationality in advanced artificial agents. \emph{Minds and Machines}, 22(2), 71--85.

\bibitem[Bostrom(2014)]{Bostrom2014} Bostrom, N. (2014). \emph{Superintelligence: Paths, Dangers, Strategies}. Oxford University Press.

\bibitem[Boyd et al.(2022)]{Boyd2022} Boyd, A.~B., Patra, A., Jarzynski, C., \& Crutchfield, J.~P. (2022). Shortcuts to thermodynamic computing. \emph{Journal of Statistical Physics}, 187, 17.

\bibitem[Brewer et al.(2011)]{Brewer2011} Brewer, J.~A., Worhunsky, P.~D., Gray, J.~R., et al. (2011). Meditation experience is associated with differences in default mode network activity and connectivity. \emph{PNAS}, 108(50), 20254--20259.

\bibitem[Bruineberg et al.(2022)]{Bruineberg2022} Bruineberg, J., Dolega, K., Dewhurst, J., \& Baltieri, M. (2022). The Emperor's new Markov blankets. \emph{Behavioral and Brain Sciences}, 45, e183.

\bibitem[Bullmore \& Sporns(2012)]{BullmoreSporns2012} Bullmore, E., \& Sporns, O. (2012). The economy of brain network organization. \emph{Nature Reviews Neuroscience}, 13, 336--349.

\bibitem[Campuzano et al.(2018)]{Campuzano2018} Campuzano, S.~A., et al. (2018). A nonlinear time series analysis of atmospheric methane mixing ratio data. \emph{Journal of Geophysical Research: Atmospheres}, 123(18), 10305--10321.

\bibitem[Cannon(2022)]{Cannon2022} Cannon, J. (2022). An enactive approach to value alignment in AI. In V.~C. Müller (Ed.), \emph{Philosophy and Theory of Artificial Intelligence 2021}. Springer.

\bibitem[Carroll-Nellenback et al.(2019)]{CarrollNellenback2019} Carroll-Nellenback, J., Frank, A., Wright, J., \& Scharf, C. (2019). The Fermi paradox and the Aurora effect: Exo-civilization settlement, expansion, and steady states. \emph{Astronomical Journal}, 158(3), 117.

\bibitem[Chen \& Fu(2023)]{ChenFu2023} Chen, X., \& Fu, F. (2023). Outlearning extortioners: Unbending strategies can foster reciprocal fairness and cooperation. \emph{PNAS Nexus}, 2(6), pgad176.

\bibitem[Cheneau et al.(2012)]{Cheneau2012} Cheneau, M., Barmettler, P., Poletti, D., et al. (2012). Light-cone-like spreading of correlations in a quantum many-body system. \emph{Nature}, 481, 484--487.

\bibitem[Ćirković(2015)]{Cirkovic2015} Ćirković, M.~M. (2015). Kardashev's classification at 50+: A fine vehicle with room for improvement. \emph{Serbian Astronomical Journal}, 191, 1--15.

\bibitem[Ćirković(2018)]{Cirkovic2018} Ćirković, M.~M. (2018). \emph{The Great Silence: Science and Philosophy of Fermi's Paradox}. Oxford University Press.

\bibitem[Clark(2013)]{Clark2013} Clark, A. (2013). Whatever next? Predictive brains, situated agents, and the future of cognitive science. \emph{Behavioral and Brain Sciences}, 36(3), 181--204.

\bibitem[Clark(2016)]{Clark2016} Clark, A. (2016). \emph{Surfing Uncertainty: Prediction, Action, and the Embodied Mind}. Oxford University Press.

\bibitem[Crawford(2021)]{Crawford2021} Crawford, K. (2021). \emph{Atlas of AI: Power, Politics, and the Planetary Costs of Artificial Intelligence}. Yale University Press.

\bibitem[Dafoe et al.(2021)]{Dafoe2021} Dafoe, A., Bachrach, Y., Hadfield, G., Horvitz, E., Larson, K., \& Graepel, T. (2021). Cooperative AI: Machines must learn to find common ground. \emph{Nature}, 593, 33--36.

\bibitem[Dick(1996)]{Dick1996} Dick, S.~J. (1996). \emph{The Biological Universe}. Cambridge University Press.

\bibitem[Donges et al.(2009)]{Donges2009} Donges, J.~F., Zou, Y., Marwan, N., \& Kurths, J. (2009). The backbone of the climate network. \emph{Europhysics Letters}, 87, 48007.

\bibitem[Ebert-Uphoff \& Deng(2012)]{EbertUphoffDeng2012} Ebert-Uphoff, I., \& Deng, Y. (2012). Causal discovery for climate research using graphical models. \emph{Journal of Climate}, 25, 5648--5665.

\bibitem[Floridi et al.(2026)]{FloridiAAIO} Floridi, L., Buttaboni, C., Gertler, N., Hine, E., Morley, J., Novelli, C., \& Schroder, T. (2026). Agentic AI Optimisation (AAIO): What it is, how it works, why it matters, and how to deal with it. \emph{Minds and Machines}, 36, 25. doi:10.1007/s11023-026-09779-8.

\bibitem[Folbre(2008)]{Folbre2008} Folbre, N. (2008). \emph{Valuing Children: Rethinking the Economics of the Family}. Harvard University Press.

\bibitem[Fonseca-Azevedo \& Herculano-Houzel(2012)]{FonsecaAzevedo2012} Fonseca-Azevedo, K., \& Herculano-Houzel, S. (2012). Metabolic constraint imposes tradeoff between body size and number of brain neurons in human evolution. \emph{PNAS}, 109(45), 18571--18576.

\bibitem[Foss-Feig et al.(2015)]{FossFeig2015} Foss-Feig, M., Gong, Z.-X., Clark, C.~W., \& Gorshkov, A.~V. (2015). Nearly-linear light cones in long-range interacting quantum systems. \emph{Physical Review Letters}, 114, 157201.

\bibitem[Fraser(2016)]{Fraser2016} Fraser, N. (2016). Contradictions of capital and care. \emph{New Left Review}, 100, 99--117.

\bibitem[Friston(2010)]{Friston2010} Friston, K. (2010). The free-energy principle: A unified brain theory? \emph{Nature Reviews Neuroscience}, 11, 127--138.

\bibitem[Froese \& Ziemke(2009)]{FroeseZiemke2009} Froese, T., \& Ziemke, T. (2009). Enactive artificial intelligence: Investigating the systemic organization of life and mind. \emph{Artificial Intelligence}, 173(3--4), 466--500.

\bibitem[Gabi et al.(2016)]{Gabi2016} Gabi, M., Neves, K., Masseron, C., et al. (2016). No relative expansion of the number of prefrontal neurons in primate and human evolution. \emph{PNAS}, 113(34), 9617--9622.

\bibitem[Gabriel(2020)]{Gabriel2020} Gabriel, I. (2020). Artificial intelligence, values, and alignment. \emph{Minds and Machines}, 30, 411--437.

\bibitem[Gallow(2025)]{Gallow2025} Gallow, J.~D. (2025). Instrumental divergence. \emph{Philosophical Studies}, 182, 1581--1607.

\bibitem[Ganesan et al.(2022)]{Ganesan2022} Ganesan, S., Beyer, E., Moffat, B., et al. (2022). Focused attention meditation in healthy adults: A systematic review and meta-analysis of cross-sectional functional MRI studies. \emph{Neuroscience \& Biobehavioral Reviews}, 141, 104846.

\bibitem[Garrison et al.(2015)]{Garrison2015} Garrison, K.~A., Zeffiro, T.~A., Scheinost, D., Constable, R.~T., \& Brewer, J.~A. (2015). Meditation leads to reduced default mode network activity beyond an active task. \emph{Cognitive, Affective, \& Behavioral Neuroscience}, 15(3), 712--720.

\bibitem[Gebru \& Torres(2024)]{GebruTorres2024} Gebru, T., \& Torres, É.~P. (2024). The TESCREAL bundle: Eugenics and the promise of utopia through artificial general intelligence. \emph{First Monday}, 29(4).

\bibitem[Godfrey-Smith(2016)]{GodfreySmith2016} Godfrey-Smith, P. (2016). \emph{Other Minds: The Octopus, the Sea, and the Deep Origins of Consciousness}. Farrar, Straus and Giroux.

\bibitem[Graeber(2011)]{Graeber2011} Graeber, D. (2011). \emph{Debt: The First 5,000 Years}. Melville House.

\bibitem[Graeber \& Wengrow(2021)]{GraeberWengrow2021} Graeber, D., \& Wengrow, D. (2021). \emph{The Dawn of Everything: A New History of Humanity}. Farrar, Straus and Giroux.

\bibitem[Hanson et al.(2021)]{Hanson2021} Hanson, R., Martin, D., McCarter, C., \& Paulson, J. (2021). If loud aliens explain human earliness, quiet aliens are also rare. \emph{Astrophysical Journal}, 922(2), 182.

\bibitem[Hauert et al.(2002)]{Hauert2002} Hauert, C., De Monte, S., Hofbauer, J., \& Sigmund, K. (2002). Volunteering as Red Queen mechanism for cooperation in public goods games. \emph{Science}, 296, 1129--1132.

\bibitem[Heims(1980)]{Heims1980} Heims, S.~J. (1980). \emph{John von Neumann and Norbert Wiener: From Mathematics to the Technologies of Life and Death}. MIT Press.

\bibitem[Hensen et al.(2015)]{Hensen2015} Hensen, B., Bernien, H., Dréau, A.~E., et al. (2015). Loophole-free Bell inequality violation using electron spins separated by 1.3 kilometres. \emph{Nature}, 526, 682--686.

\bibitem[Herculano-Houzel(2011)]{HerculanoHouzel2011} Herculano-Houzel, S. (2011). Scaling of brain metabolism with a fixed energy budget per neuron. \emph{PLoS ONE}, 6(3), e17514.

\bibitem[Herculano-Houzel(2012)]{HerculanoHouzel2012} Herculano-Houzel, S. (2012). The remarkable, yet not extraordinary, human brain as a scaled-up primate brain and its associated cost. \emph{PNAS}, 109(Suppl 1), 10661--10668.

\bibitem[Herculano-Houzel et al.(2010)]{HerculanoHouzel2010} Herculano-Houzel, S., Mota, B., Wong, P., \& Kaas, J.~H. (2010). Connectivity-driven white matter scaling and folding in primate cerebral cortex. \emph{PNAS}, 107(44), 19008--19013.

\bibitem[Hilbe et al.(2013)]{Hilbe2013} Hilbe, C., Nowak, M.~A., \& Sigmund, K. (2013). Evolution of extortion in iterated prisoner's dilemma games. \emph{PNAS}, 110(17), 6913--6918.

\bibitem[Hilbe et al.(2014a)]{Hilbe2014a} Hilbe, C., Wu, B., Traulsen, A., \& Nowak, M.~A. (2014). Cooperation and control in multiplayer social dilemmas. \emph{PNAS}, 111(46), 16425--16430.

\bibitem[Hilbe et al.(2014b)]{Hilbe2014b} Hilbe, C., Röhl, T., \& Milinski, M. (2014). Extortion subdues human players but is finally punished in the prisoner's dilemma. \emph{Nature Communications}, 5, 3976.

\bibitem[Hofman(2014)]{Hofman2014} Hofman, M.~A. (2014). Evolution of the human brain: When bigger is better. \emph{Frontiers in Neuroanatomy}, 8, 15.

\bibitem[Hohwy(2013)]{Hohwy2013} Hohwy, J. (2013). \emph{The Predictive Mind}. Oxford University Press.

\bibitem[Hubinger et al.(2019)]{Hubinger2019} Hubinger, E., van Merwijk, C., Mikulik, V., Skalse, J., \& Garrabrant, S. (2019). Risks from learned optimization in advanced machine learning systems. arXiv:1906.01820.

\bibitem[Hyvönen et al.(2018)]{Hyvonen2018} Hyvönen, S., et al. (2018). Wavelet-based mutual information analysis of atmospheric surface-layer fluxes. \emph{Boundary-Layer Meteorology}, 167(3), 345--363.

\bibitem[Jackson \& Criado-Perez(2024)]{JacksonCriadoPerez2024} Jackson, C.~J., \& Criado-Perez, C. (2024). Why the Fermi paradox may not be well explained by Wong and Bartlett's theory of civilization collapse. \emph{Journal of the Royal Society Interface}, 21(219), 20240140.

\bibitem[Kardashev(1964)]{Kardashev1964} Kardashev, N.~S. (1964). Transmission of information by extraterrestrial civilizations. \emph{Soviet Astronomy}, 8, 217--221.

\bibitem[Karst et al.(2023)]{Karst2023} Karst, J., Jones, M.~D., \& Hoeksema, J.~D. (2023). Positive citation bias and overinterpreted results lead to misinformation on common mycorrhizal networks in forests. \emph{Nature Ecology \& Evolution}, 7, 501--511.

\bibitem[Kraskov et al.(2004)]{Kraskov2004} Kraskov, A., Stögbauer, H., \& Grassberger, P. (2004). Estimating mutual information. \emph{Physical Review E}, 69, 066138.

\bibitem[Lamport(1978)]{Lamport1978} Lamport, L. (1978). Time, clocks, and the ordering of events in a distributed system. \emph{Communications of the ACM}, 21(7), 558--565.

\bibitem[Landauer(1961)]{Landauer1961} Landauer, R. (1961). Irreversibility and heat generation in the computing process. \emph{IBM Journal of Research and Development}, 5(3), 183--191.

\bibitem[Levin(2023)]{Levin2023} Levin, M. (2023). Bioelectric networks: The cognitive glue enabling evolutionary scaling from physiology to mind. \emph{Animal Cognition}, 26, 1865--1891.

\bibitem[Lieb \& Robinson(1972)]{LiebRobinson1972} Lieb, E.~H., \& Robinson, D.~W. (1972). The finite group velocity of quantum spin systems. \emph{Communications in Mathematical Physics}, 28(3), 251--257.

\bibitem[Luppi et al.(2024)]{Luppi2024} Luppi, A.~I., Mediano, P.~A.~M., Rosas, F.~E., et al. (2024). A synergistic workspace for human consciousness revealed by integrated information decomposition. \emph{eLife}, 12, RP88173.

\bibitem[Lyon et al.(2021)]{Lyon2021} Lyon, P., Keijzer, F., Arendt, D., \& Levin, M. (2021). Reframing cognition: Getting down to biological basics. \emph{Philosophical Transactions of the Royal Society B}, 376(1820), 20190750.

\bibitem[Man \& Damasio(2019)]{ManDamasio2019} Man, K., \& Damasio, A. (2019). Homeostasis and soft robotics in the design of feeling machines. \emph{Nature Machine Intelligence}, 1(10), 446--452.

\bibitem[Mauss(1925)]{Mauss1925} Mauss, M. (1925). \emph{Essai sur le don}. L'Année Sociologique.

\bibitem[McCloskey(2010)]{McCloskey2010} McCloskey, D.~N. (2010). \emph{Bourgeois Dignity}. University of Chicago Press.

\bibitem[McCloskey(2016)]{McCloskey2016} McCloskey, D.~N. (2016). \emph{Bourgeois Equality}. University of Chicago Press.

\bibitem[McMillen \& Levin(2024)]{McMillenLevin2024} McMillen, P., \& Levin, M. (2024). Collective intelligence: A unifying concept for integrating biology across scales and substrates. \emph{Communications Biology}, 7, 378.

\bibitem[Mini et al.(2023)]{Mini2023} Mini, U., Grietzer, P., Sharma, M., Meek, A., MacDiarmid, M., \& Turner, A.~M. (2023). Understanding and controlling a maze-solving policy network. arXiv:2310.08043.

\bibitem[Mirowski(2002)]{Mirowski2002} Mirowski, P. (2002). \emph{Machine Dreams: Economics Becomes a Cyborg Science}. Cambridge University Press.

\bibitem[Mokyr(2016)]{Mokyr2016} Mokyr, J. (2016). \emph{A Culture of Growth: The Origins of the Modern Economy}. Princeton University Press.

\bibitem[Mota \& Herculano-Houzel(2015)]{MotaHerculanoHouzel2015} Mota, B., \& Herculano-Houzel, S. (2015). Cortical folding scales universally with surface area and thickness. \emph{Science}, 349(6243), 74--77.

\bibitem[Müller \& Cannon(2022)]{MullerCannon2022} Müller, V.~C., \& Cannon, M. (2022). Existential risk from AI and orthogonality: Can we have it both ways? \emph{Ratio}, 35(1), 25--36.

\bibitem[Ngo(2019)]{Ngo2019} Ngo, R. (2019). Coherent behaviour in the real world is an incoherent concept. \emph{AI Alignment Forum}.

\bibitem[Nowak(2006)]{Nowak2006} Nowak, M.~A. (2006). Five rules for the evolution of cooperation. \emph{Science}, 314, 1560--1563.

\bibitem[Nowak \& Sigmund(1993)]{NowakSigmund1993} Nowak, M., \& Sigmund, K. (1993). A strategy of win-stay, lose-shift that outperforms tit-for-tat in the prisoner's dilemma. \emph{Nature}, 364, 56--58.

\bibitem[Omohundro(2008)]{Omohundro2008} Omohundro, S.~M. (2008). The basic AI drives. In P. Wang, B. Goertzel, \& S. Franklin (Eds.), \emph{Proceedings of the First AGI Conference}, 483--492.

\bibitem[Ostrom(1990)]{Ostrom1990} Ostrom, E. (1990). \emph{Governing the Commons}. Cambridge University Press.

\bibitem[Pearl(1988)]{Pearl1988} Pearl, J. (1988). \emph{Probabilistic Reasoning in Intelligent Systems: Networks of Plausible Inference}. Morgan Kaufmann.

\bibitem[Pihlakas \& Pyykkö(2024)]{PihlakasPyykko2024} Pihlakas, R., \& Pyykkö, J. (2024). From homeostasis to resource sharing: Biologically and economically aligned multi-objective multi-agent AI safety benchmarks. arXiv:2410.00081.

\bibitem[Polanyi(1944)]{Polanyi1944} Polanyi, K. (1944). \emph{The Great Transformation}. Beacon Press.

\bibitem[Press \& Dyson(2012)]{PressDyson2012} Press, W.~H., \& Dyson, F.~J. (2012). Iterated prisoner's dilemma contains strategies that dominate any evolutionary opponent. \emph{PNAS}, 109(26), 10409--10413.

\bibitem[Rahrig et al.(2022)]{Rahrig2022} Rahrig, H., Vago, D.~R., Passarelli, M.~A., et al. (2022). Meta-analytic evidence that mindfulness training alters resting state default mode network connectivity. \emph{Scientific Reports}, 12, 12260.

\bibitem[Reid(2023)]{Reid2023} Reid, C.~R. (2023). Thoughts from the forest floor: A review of cognition in the slime mould \emph{Physarum polycephalum}. \emph{Animal Cognition}, 26(6), 1783--1797.

\bibitem[Rieder(2008)]{Rieder2008} Rieder, J. (2008). \emph{Colonialism and the Emergence of Science Fiction}. Wesleyan University Press.

\bibitem[Robinson et al.(2024)]{Robinson2024} Robinson, D.~G., Ammer, C., Polle, A., et al. (2024). Mother trees, altruistic fungi, and the perils of plant personification. \emph{Trends in Plant Science}, 29(1), 20--31.

\bibitem[Runge et al.(2015)]{Runge2015} Runge, J., Petoukhov, V., Donges, J.~F., et al. (2015). Identifying causal gateways and mediators in complex spatio-temporal systems. \emph{Nature Communications}, 6, 8502.

\bibitem[Sagan \& Newman(1983)]{SaganNewman1983} Sagan, C., \& Newman, W.~I. (1983). The solipsist approach to extraterrestrial intelligence. \emph{Quarterly Journal of the Royal Astronomical Society}, 24, 113--121.

\bibitem[Sandberg et al.(2018)]{Sandberg2018} Sandberg, A., Drexler, E., \& Ord, T. (2018). Dissolving the Fermi paradox. arXiv:1806.02404.

\bibitem[Schwieterman et al.(2024)]{Schwieterman2024} Schwieterman, E.~W., Haqq-Misra, J., Kopparapu, R.~K., et al. (2024). Artificial greenhouse gases as exoplanet technosignatures. \emph{Astrophysical Journal}, 969(1), 20.

\bibitem[Seifert et al.(2024)]{Seifert2024} Seifert, G., Sealander, A., Marzen, S., \& Levin, M. (2024). From reinforcement learning to agency: Frameworks for understanding basal cognition. \emph{BioSystems}, 235, 105107.

\bibitem[Sen(1977)]{Sen1977} Sen, A. (1977). Rational fools: A critique of the behavioral foundations of economic theory. \emph{Philosophy \& Public Affairs}, 6(4), 317--344.

\bibitem[Shah(2018)]{Shah2018} Shah, R. (2018). Coherence arguments do not imply goal-directed behavior. \emph{AI Alignment Forum}.

\bibitem[Sharadin(2025)]{Sharadin2025} Sharadin, N. (2025). Promotionalism, orthogonality, and instrumental convergence. \emph{Philosophical Studies}, 182(7), 1725--1755.

\bibitem[Simard et al.(2025)]{Simard2025} Simard, S.~W., Ryan, T.~L., \& Perry, D.~A. (2025). Opinion: Response to questions about common mycorrhizal networks. \emph{Frontiers in Forests and Global Change}, 7, 1512518.

\bibitem[Sivak \& Crooks(2012)]{SivakCrooks2012} Sivak, D.~A., \& Crooks, G.~E. (2012). Thermodynamic metrics and optimal paths. \emph{Physical Review Letters}, 108, 190602.

\bibitem[Sterling(2012)]{Sterling2012} Sterling, P. (2012). Allostasis: A model of predictive regulation. \emph{Physiology \& Behavior}, 106(1), 5--15.

\bibitem[Stewart \& Plotkin(2013)]{StewartPlotkin2013} Stewart, A.~J., \& Plotkin, J.~B. (2013). From extortion to generosity: Evolution in the iterated prisoner's dilemma. \emph{PNAS}, 110(38), 15348--15353.

\bibitem[Tarsney(2025)]{Tarsney2025} Tarsney, C. (2025). Will artificial agents pursue power by default? arXiv:2506.06352.

\bibitem[Tero et al.(2010)]{Tero2010} Tero, A., Takagi, S., Saigusa, T., et al. (2010). Rules for biologically inspired adaptive network design. \emph{Science}, 327, 439--442.

\bibitem[Thompson(2007)]{Thompson2007} Thompson, E. (2007). \emph{Mind in Life}. Harvard University Press.

\bibitem[Thorstad(2024)]{Thorstad2024} Thorstad, D. (2024). What power-seeking theorems do not show. \emph{Global Priorities Institute Working Paper}.

\bibitem[Tipler(1980)]{Tipler1980} Tipler, F.~J. (1980). Extraterrestrial intelligent beings do not exist. \emph{Quarterly Journal of the Royal Astronomical Society}, 21, 267--281.

\bibitem[Tran et al.(2021)]{Tran2021} Tran, M.~C., Guo, A.~Y., Su, Y., et al. (2021). Locality and digital quantum simulation of power-law interactions. \emph{Physical Review Letters}, 127, 160401.

\bibitem[Tronto(1993)]{Tronto1993} Tronto, J.~C. (1993). \emph{Moral Boundaries: A Political Argument for an Ethic of Care}. Routledge.

\bibitem[Turner et al.(2021)]{Turner2021} Turner, A.~M., Smith, L., Shah, R., Critch, A., \& Tadepalli, P. (2021). Optimal policies tend to seek power. \emph{Advances in Neural Information Processing Systems}, 34. arXiv:1912.01683.

\bibitem[Turner \& Tadepalli(2022)]{TurnerTadepalli2022} Turner, A.~M., \& Tadepalli, P. (2022). Parametrically retargetable decision-makers tend to seek power. \emph{Advances in Neural Information Processing Systems}, 35.

\bibitem[Turner(2022a)]{Turner2022a} Turner, A.~M. (2022a). When most VNM-coherent preference orderings have convergent instrumental incentives. \emph{AI Alignment Forum}.

\bibitem[Turner(2022b)]{Turner2022b} Turner, A.~M. (2022b). Reward is not the optimization target. \emph{AI Alignment Forum}.

\bibitem[Turner(2022c)]{Turner2022c} Turner, A.~M. (2022c). A certain formalization of corrigibility is VNM-incoherent. \emph{AI Alignment Forum}.

\bibitem[Turner et al.(2023)]{Turner2023} Turner, A.~M., Thiergart, L., Leech, G., et al. (2023). Activation addition: Steering language models without optimization. arXiv:2308.10248.

\bibitem[Van Dam et al.(2018)]{VanDam2018} Van Dam, N.~T., van Vugt, M.~K., Vago, D.~R., et al. (2018). Mind the hype: A critical evaluation and prescriptive agenda for research on mindfulness and meditation. \emph{Perspectives on Psychological Science}, 13(1), 36--61.

\bibitem[Varian(2005)]{Varian2005} Varian, H.~R. (2005). Revealed preference. In M. Szenberg et al. (Eds.), \emph{Samuelsonian Economics and the Twenty-First Century}. Oxford University Press.

\bibitem[Watson \& Levin(2023)]{WatsonLevin2023} Watson, R., \& Levin, M. (2023). The collective intelligence of evolution and development. \emph{Collective Intelligence}, 2(2), 1--22.

\bibitem[Wong \& Bartlett(2022)]{WongBartlett2022} Wong, M.~L., \& Bartlett, S. (2022). Asymptotic burnout and homeostatic awakening: A possible solution to the Fermi paradox? \emph{Journal of the Royal Society Interface}, 19, 20220029.

\bibitem[Wright(2020)]{Wright2020} Wright, J.~T. (2020). Dyson spheres. \emph{Serbian Astronomical Journal}, 200, 1--18.

\bibitem[Wright(2023)]{Wright2023} Wright, J.~T. (2023). Application of the thermodynamics of radiation to Dyson spheres. \emph{Astrophysical Journal}, 956, 34.

\end{thebibliography}

\end{document}
